 \documentclass[11pt,letterpaper,usenames,dvipsnames]{article}           
%   %%%%%%%%%%%%%%%%%%%%%%%%%%%%%%%%%%%%%%%%%%%%%%%%%%%%%%%%%%%%%%%%%%%%%%%%%%%%%%%
% PACKAGES
\usepackage{amsmath,amssymb,amsthm}
\allowdisplaybreaks
\usepackage{aliascnt}
\usepackage{fullpage}
\usepackage{multirow}
\usepackage{bbm,soul}
\usepackage[usenames,dvipsnames]{xcolor}
\usepackage{lipsum}
\usepackage{makecell}
\usepackage{bm}
\usepackage{braket}
\usepackage{adjustbox,mdframed}
\usepackage[operators,sets,primitives,asymptotics,lambda,keys,ff]{cryptocode}
\usepackage{enumitem}
\usepackage{geometry}
\usepackage{booktabs}
\usepackage{siunitx}
\usepackage[font=footnotesize,labelfont=bf]{caption}
\usepackage{subcaption}
\usepackage{mleftright}
\usepackage{wrapfig}
\usepackage{setspace}
\usepackage{xspace}
\usepackage{algorithm,algorithmic}
\floatname{algorithm}{Protocol}
%\crefname{algorithm}{Protocol}{Protocols}
%\Crefname{algorithm}{Protocol}{Protocols}
\renewcommand{\algorithmiccomment}[1]{// #1}
\renewcommand{\algorithmicrequire}{\textbf{Input:}}
\renewcommand{\algorithmicensure}{\textbf{Output:}}
\usepackage{authblk}

\usepackage{framed}
\usepackage{comment}
\usepackage{placeins} % in the preamble

\usepackage{float}
\usepackage{graphicx}
\graphicspath{ {./figures/} }
\usepackage[T1]{fontenc}
\usepackage{titlesec}
\usepackage{stmaryrd}

\usepackage[backend=biber,giveninits,style=alphabetic,maxalphanames=6,maxnames=6,backref]{biblatex}
\addbibresource{references.bib} 
%\addbibresource{cryptobib/abbrev3.bib} 
%\addbibresource{cryptobib/crypto_crossref.bib}

%%%%%%%%%%%%%%%%%%%%%%%%%%%%%%%%%%%%%%%%%%%%%%%%%%%%%%%%%%%%%%%%%%%%%%%%%%%%%%%%

\newcommand{\plaintitle}{$\FF 2\ZZ$}

\PassOptionsToPackage{hyphens}{url}\usepackage{hyperref}
\hypersetup{
  colorlinks=true,
  linkcolor=PineGreen,%RedViolet,
  breaklinks=true,
  citecolor=PineGreen,      
  urlcolor=PineGreen,%RedViolet,
  pdftitle={\plaintitle},
  pdfpagemode=FullScreen,
}
\DefineBibliographyStrings{english}{%
backrefpage = {cited on p.},% originally "cited on page"
backrefpages = {cited on pp.},% originally "cited on pages"
}
% Remove starred sections from toc
\DeclareRobustCommand{\SkipTocEntry}[5]{} 
\usepackage[capitalize, nameinlink]{cleveref}
\usepackage{tikz}
\usetikzlibrary{positioning, arrows.meta, calc}

\DeclareMathOperator{\im}{im}

\newcommand{\psiv}[1]{{\color{Emerald}[#1 -- Psi.]}}
\newcommand{\Albert}[1]{{\color{purple}Albert: #1}}
\newcommand{\ulrich}[1]{{\color{brown}Ulrich: #1}}
\newcommand{\osdnk}[1]{{\color{red}@osdnk: #1}}
\newcommand{\kai}[1]{{\color{ForestGreen}Kai: #1}}
\newcommand{\Aru}[1]{{\color{blue}Aru: #1}}
\newcommand{\remco}[1]{{\color{orange}Remco: #1}}
\newcommand{\john}[1]{{\color{pink}John: #1}}

% Shorthand for \mathcal
\newcommand{\mc}[1]{\mathcal{#1}}

\newcommand{\gp}{\mathsf{gp}}
\newcommand{\commit}{\mathsf{Commit}}
\newcommand{\open}{\mathsf{Open}}
\newcommand{\evaluation}{\mathsf{Evaluation}}
\newcommand{\testingP}{\mathsf{Testing Phase}}
\newcommand{\evaluationP}{\mathsf{Evaluation Phase}}
\newcommand{\com}{\mathsf{com}}

\newcommand{\ad}{\mathsf{AD}}

\newcommand{\code}{\mathcal{C}}

\newcommand{\polyspace}{(\QQ[X])[\vec Y]}


\newcommand{\encode}{\mathsf{Enc}}


\newcommand{\size}{2^\mu}
\newcommand{\matrixU}{\mathbf{\hat{u}}}
\newcommand{\matrixV}{\mathbf{V}}
\newcommand{\vectorUi}{\mathbf{\hat{u}_i}}
\newcommand{\oracleUi}{\llbracket\mathbf{\hat{u}}_i\rrbracket}

%%%%%%%%%%%%%%%%%%%%%%%%%%%%%%%%%%%%%%%%%%%%%%%%%%%%%%%%%%%%%%%%%%%%%%%%%%%%%%%%
% FORMATTING
%

\interfootnotelinepenalty=10000

\newcommand{\keywords}[1]{\bigskip\par\noindent{\textbf{Keywords\/}: #1}}

%%%%%%%%%%%%%%%%%%%%%%%%%%%%%%%%%%%%%%%%%%%%%%%%%%%%%%%%%%%%%%%%%%%%%%%%%%%%%%%%
% THEOREMS
\theoremstyle{plain}
\newtheorem{theorem}{Theorem}[section]

\newaliascnt{lemma}{theorem}
\newtheorem{lemma}[lemma]{Lemma}
\aliascntresetthe{lemma}
\crefname{lemma}{Lemma}{Lemmas}
\Crefname{lemma}{Lemma}{Lemmas}

\newaliascnt{proposition}{theorem}
\newtheorem{proposition}[proposition]{Proposition}
\aliascntresetthe{proposition}
\crefname{proposition}{Proposition}{Propositions}
\Crefname{proposition}{Proposition}{Propositions}

\newaliascnt{corollary}{theorem}
\newtheorem{corollary}[corollary]{Corollary}
\aliascntresetthe{corollary}
\crefname{corollary}{Corollary}{Corollaries}
\Crefname{corollary}{Corollary}{Corollaries}

\newaliascnt{claim}{theorem}
\newtheorem{claim}[claim]{Claim}
\aliascntresetthe{claim}
\crefname{claim}{Claim}{Claims}
\Crefname{claim}{Claim}{Claims}

\newaliascnt{conjecture}{theorem}
\newtheorem{conjecture}[conjecture]{Conjecture}
\aliascntresetthe{conjecture}
\crefname{conjecture}{Conjecture}{Conjectures}
\Crefname{conjecture}{Conjecture}{Conjectures}

\newaliascnt{question}{theorem}
\newtheorem{question}[question]{Question}
\aliascntresetthe{question}
\crefname{question}{Question}{Questions}
\Crefname{question}{Question}{Questions}


\theoremstyle{definition}
\newaliascnt{definition}{theorem}
\newtheorem{definition}[definition]{Definition}
\aliascntresetthe{definition}
\crefname{definition}{Definition}{Definitions}
\Crefname{definition}{Definition}{Definitions}

\theoremstyle{definition}
\newaliascnt{convention}{theorem}
\newtheorem{convention}[convention]{Convention}
\aliascntresetthe{convention}
\crefname{convention}{Convention}{Conventions}
\Crefname{convention}{Convention}{Conventions}

\theoremstyle{definition}
\newaliascnt{construction}{theorem}
\newtheorem{construction}[construction]{Construction}
\aliascntresetthe{construction}
\crefname{construction}{Construction}{Constructions}
\Crefname{construction}{Construction}{Constructions}

\theoremstyle{definition}
\newaliascnt{remark}{theorem}
\newtheorem{remark}[remark]{Remark}
\aliascntresetthe{remark}
\crefname{remark}{Remark}{Remarks}
\Crefname{remark}{Remark}{Remarks}

\newtheorem{fact}[theorem]{Fact}

%%%%%%%%%%%%%%%%%%%%%%%%%%%%%%%%%%%%%%%%%%%%%%%%%%%%%%%%%%%%%%%%%%%%%%%%%%%%%%%%
% BREAKABLE ALGORITHM ENVIRONMENT
% Allows algorithms/protocols to span multiple pages
\makeatletter
\newenvironment{breakablealgorithm}
  {\begin{center}
     \refstepcounter{algorithm}%
     \hrule height.8pt depth0pt \kern2pt
     \renewcommand{\caption}[2][\relax]{%
       {\raggedright\textbf{\ALG@name~\thealgorithm} ##2\par}%
       \ifx\relax##1\relax
         \addcontentsline{loa}{algorithm}{\protect\numberline{\thealgorithm}##2}%
       \else
         \addcontentsline{loa}{algorithm}{\protect\numberline{\thealgorithm}##1}%
       \fi
       \kern2pt\hrule\kern2pt
     }
  }{\kern2pt\hrule\relax\end{center}}
\makeatother

\newaliascnt{example}{theorem}
\newtheorem{example}[example]{Example}
\aliascntresetthe{example}
\crefname{example}{Example}{Examples}
\Crefname{example}{Example}{Examples}
%%%%%%%%%%%%%%%%%%%%%%%%%%%%%%%%%%%%%%%%%%%%%%%%%%%%%%%%%%%%%%%%%%%%%%%%%%%%%%%%
% Make \emph meaningful in italic environments (e.g., when using \theoremstyle{plain})
\renewcommand\eminnershape{\itshape\sffamily}

\newcommand{\emath}[1]{\ensuremath{#1}}
\newcommand{\newc}[2]{\newcommand{#1}{\ensuremath{#2}\xspace}}
\newcommand{\renewc}[2]{\renewcommand{#1}{\ensuremath{#2}\xspace}}

% allow in-line math breaks
\newcommand{\CB}{\allowbreak}

% such that in probability statements
\newcommand{\pST}{\; \middle| \;}
\DeclareMathOperator{\Concat}{\; || \; }
% \DeclareMathOperator{\max}{max}
\newcommand{\Or}{\vee}
\renewcommand{\and}{\ \mathsf{AND} \ }
\newcommand{\Not}{\neg}
\newc{\defeq}{:=}
\newcommand{\Union}{\cup}
\newcommand{\defemph}[1]{\textbf{\emph{#1}}}

% some math macros
\newcommand{\List}[1]{\left[#1\right]}
\newcommand{\SmallList}[1]{[#1]}
\renewcommand{\set}[1]{\ensuremath{\mleft\{#1\mright\}}}
\newcommand{\Range}[2] {[#1 \, .. \, #2]}
\newcommand{\OpenOpen}[2]{\left[#1, #2\right]}
\newcommand{\OpenClosed}[2]{\left[#1, #2\right)}
\newcommand{\ClosedOpen}[2]{\left(#1, #2\right]}
\newcommand{\ClosedClosed}[2]{\left(#1, #2\right)}

\newcommand{\Element}[3]{#1_{#2,#3}}

% tuples with in-line math breaks
\newcommand{\tuple}[2]{(#1 ,\CB \dots,\CB #2)}
\newcommand{\pair}[2]{(#1 ,\CB #2)}
\newcommand{\ppair}[2]{\left(#1 ,\CB #2\right)}
\newcommand{\triple}[3]{(#1 ,\CB #2,\CB #3)}
\newcommand{\ptriple}[3]{\left(#1 ,\CB #2,\CB #3\right)}
\newcommand{\quadruple}[4]{(#1 ,\CB #2,\CB #3,\CB #4)}
\newcommand{\pquadruple}[4]{\left(#1 ,\CB #2,\CB #3,\CB #4\right)}
\newcommand{\quintuple}[5]{(#1 ,\CB #2,\CB #3,\CB #4,\CB #5)}
\newcommand{\quintuplenp}[5]{#1 ,\CB #2,\CB #3,\CB #4,\CB #5}
\newcommand{\sextuple}[6]{(#1 ,\CB #2,\CB #3,\CB #4,\CB #5,\CB #6)}
\newcommand{\sextuplenp}[6]{#1 ,\CB #2,\CB #3,\CB #4,\CB #5,\CB #6}
\newcommand{\septuple}[7]{(#1 ,\CB #2,\CB #3,\CB #4,\CB #5,\CB #6,\CB #7)}
\newcommand{\septuplenp}[7]{#1 ,\CB #2,\CB #3,\CB #4,\CB #5,\CB #6,\CB #7}
\newcommand{\octuple}[8]{(#1 ,\CB #2,\CB #3,\CB #4,\CB #5,\CB #6,\CB #7,\CB #8)}
\newcommand{\octuplenp}[8]{#1 ,\CB #2,\CB #3,\CB #4,\CB #5,\CB #6,\CB #7,\CB #8}
\newcommand{\nonuple}[9]{(#1 ,\CB #2,\CB #3,\CB #4,\CB #5,\CB #6,\CB #7,\CB #8,\CB #9)}
\newcommand{\nonuplenp}[9]{#1 ,\CB #2,\CB #3,\CB #4,\CB #5,\CB #6,\CB #7,\CB #8,\CB #9}

% expectation
\newcommand{\prob}[1]{\Pr\left[ #1 \right]}
\newcommand{\Expectation}{\mathbb{E}}

% random variable
\newcommand{\RandomVariable}{X}

% complexity classes
\newcommand{\Class}[1]{\mathsf{#1}}
\newcommand{\NP}{\Class{NP}}
\newcommand{\DTIME}[1]{\Class{DTIME}(#1)}

% types
\newc{\bits}{\{0,1\}}
\newc{\strings}{\bits^{*}}
\newc{\naturals}{\mathbb{N}}
\newc{\cmplx}{\mathbb{C}}
\newc{\reals}{\mathbb{R}}
\newc{\rationals}{\mathbb{Q}}
\newc{\integers}{\mathbb{Z}}
\newc{\integersq}{\mathbb{Z}_{\modulus}}
\newc{\integersp}{\mathbb{Z}_{\pmodulus}}
\newc{\integersalpha}{\mathbb{Z}_{\alpha}}
\newcommand{\integerball}[1]{\integers(#1)}
\newc{\preimagepowers}{P}

\newcommand{\EmptyString}{\varepsilon}

\newcommand{\field}{\mathbb{F}}
\newc{\extend}{\mathsf{Extend}}
\newc{\ff}{\mathbb{F}_{\pmodulus}}
\newcommand{\Pairing}[2]{\PairingSymb(#1,\ #2)}
\newcommand{\fixMSM}{\textrm{f-MSM}}
\newcommand{\varMSM}{\textrm{v-MSM}}
\newcommand{\FFT}{\mathsf{FFT}}
\newcommand{\IFFT}{\textrm{IFFT}}

\newcommand{\Max}{\mathsf{Max}}

\newcommand{\innerproduct}[2]{\langle #1, \ #2 \rangle}

\newc{\randpick}{\ {\gets}{\scriptstyle\$}\ }

\newcommand{\Size}[1]{\abs{#1}}
\newcommand{\Time}{t}

\renewc{\secparam}{\lambda}
\newc{\soundnesserror}{\epsilon}
\newc{\knowledgeerror}{\kappa}
\newc{\extractiontime}{\mathsf{ext}}
\newc{\rbrextractor}{\mathcal{E}_{\mathsf{rbr}}}
\newc{\kstate}{\mathsf{KState}}
\newc{\kswit}{\mathsf{w}}
\newc{\provermsg}{\pi}
\newc{\verifiermsg}{\rho}
\newc{\provertime}{\mathsf{pt}}
\newc{\verifiertime}{\mathsf{vt}}
\newcommand{\Deg}{\mathrm{deg}}
\newcommand{\VIEW}{\mathrm{View}}

% attackers
\newcommand{\QuerySampler}{\mathcal{Q}}
\newc{\adversary}{\mathcal{A}}
\newcommand{\adversaryb}{\mathcal{B}}
\newcommand{\adversaryd}{\mathcal{D}}
\newc{\experiment}{\mathsf{Exp}}
\newcommand{\Simulator}{\mathcal{S}}
\newcommand{\PCSimulator}{\mathcal{S}_{\PC}}
\newcommand{\SPCSimulator}{\mathcal{S}_{\mathsf{s}}}
\newcommand{\SqPCSimulator}{\mathcal{S}_{\mathsf{m}}}
\newcommand{\Real}{\mathsf{Real}}
\newcommand{\Ideal}{\mathcal{I}}
\newc{\prim}{\mathsf{Prim}}
\newc{\generator}{\pcvarstyle{j}}
\newcommand{\principalideal}[1]{\langle \generator_{#1} \rangle}
\newcommand{\primeideal}{\mathfrak{I}}
\newc{\extractoradversary}{\mathcal{E}_{\adversary}}
\newcommand{\Challenger}{\mathcal{C}}

% oracles
\newcommand{\HashOracle}{\rho}
\newcommand{\Oracle}{\mathcal{O}}
\newcommand{\OracleFamily}{\mathbb{O}}
\newcommand{\SuccessOracles}{\mathsf{S}}

% public parameters
\newcommand{\CRS}{\mathsf{crs}}
\newcommand{\td}{\mathsf{td}}
\newcommand{\PK}{\mathsf{PK}}

% relation, instance, witness
\newcommand{\Relation}{\mathcal{R}}
\newcommand{\Language}{\mathcal{L}}
\newcommand{\Index}{\mathbbm{i}}
\newcommand{\Instance}{\mathbbm{x}}
\newcommand{\InstanceSizeBound}{\mathsf{N}}
\newcommand{\RelationWithSizeBound}{\Relation_{\InstanceSizeBound}}

\newcommand{\ValidInstance}{\Language(\Relation)}
\newcommand{\ValidInstanceWithSizeBound}{\Language(\RelationWithSizeBound)}

\newcommand{\AuxiliaryInp}{\mathbbm{z}}
\newcommand{\AuxiliaryInpDistri}{\mathcal{Z}}
\newcommand{\aux}{\mathsf{aux}}


\newcommand{\polylog}{\mathrm{polylog}}
\newcommand{\hybrid}[1]{\ensuremath{\mathsf{Hyb}_{#1}\xspace}}

\newcommand{\vhl}[1]{{\color{blue} #1}}
\newcommand{\mhl}[1]{\text{\hl{$#1$}}}
\newcommand{\iseq}{\overset{?}{=}}
\newcommand{\isleq}{\overset{?}{\leq}}

% Hash families
\newcommand{\HashFamily}{\mathcal{H}}
\newcommand{\HashIndex}{I}
\newcommand{\HashKey}{I}
\newcommand{\HashIndexSet}{\mathcal{I}}
\newcommand{\Hash}{h}
\newcommand{\HashKeySpace}{\mathcal{I}}

% Transcripts
\newcommand{\TranscriptSet}{\mathcal{X}}
\newcommand{\challengeset}{\mathcal{H}}
\newcommand{\State}{\mathsf{state}}
\newcommand{\Transcript}{\tau}
\newcommand{\Tree}{\textnormal{\textsc{tree}}}
\newcommand{\EmptyTranscript}{\emptyset}
\newcommand{\StatesMultiRounds}[1]{[\State_i]_{i=1}^{\PCRound}}
\newcommand{\QueryTranscript}{\mathsf{qt}}
\newcommand{\Query}{q}
\newcommand{\NumQueries}{Q}
\newcommand{\QuerySet}{\mathcal{Q}}
\newcommand{\ChallengeSetSize}{p}
\newcommand{\Image}{\mathsf{Image}}

% Algebras
\newc{\ring}{R}
\newc{\ideal}{\mathfrak{n}}
\newc{\ringdegree}{n}
\newc{\modulus}{q}
\newc{\pmodulus}{M}
\newc{\modq}{\bmod\modulus}
\newc{\modp}{\bmod\pmodulus}
\newc{\fring}{\ring_{\modulus}}
\newc{\pring}{\ring_{\pmodulus}}
\newc{\fringinvertiblesubset}{\ring_{\modulus}^{\times}}
\newc{\fringmultgroup}{\fring^{\times}}
\newc{\embedding}{\iota}
\newcommand{\automorphism}[1]{\emath{\sigma_{-1}(#1)}}
\newcommand{\ltwonorm}[1]{\norm{#1}_{2}}
\newcommand{\linfinity}{\ell_\infty}
\newcommand{\linfinitynorm}[1]{\norm{#1}_{\infty}}

% Codes
\renewc{\dist}{\mathsf{dist}}
\renewc{\dim}{\mathsf{dim}}
\newc{\blocklength}{\mathsf{n}}
\newcommand{\reldist}[2]{\Delta\left(#1,#2\right)}

% Footnotes
\newcommand{\astfootnote}[1]{%
\let\oldthefootnote=\thefootnote%
\setcounter{footnote}{0}%
\renewcommand{\thefootnote}{\fnsymbol{footnote}}%
\footnote{#1}%
\addtocounter{footnote}{-1}%
\let\thefootnote=\oldthefootnote%
}

\newcommand\blfootnote[1]{%
  \begingroup
  \renewcommand\thefootnote{}\footnote{#1}%
  \addtocounter{footnote}{-1}%
  \endgroup
}

% Better paragraphs headers
\newcommand{\parhead}[1]{\medskip\noindent\textbf{#1}.\quad}

\newcommand{\newdefs}[1] {\setlength{\fboxsep}{1pt}\colorbox{gray!20}{\(#1\)}}


%general formatting
\newcommand{\pcvarstyle}[1]{\mathsf{#1}}
\newcommand{\continue}{{\Huge{\hl{$\cdots$}}}}
% \renewcommand{\case}[1]{\medskip\noindent{\fbox{Case #1:}}}
% \newcommand{\ncase}[1]{\medskip\noindent{\fbox{#1:}}}
% \newcommand{\ngame}[1]{\medskip\noindent{\fbox{Game $\game{#1}$:}}}
\newcommand{\mhyph}{\text{-}}
\newcommand{\ncase}[1]{\medskip\noindent{\textbf{#1:}}}
\newcommand{\ngame}[1]{\medskip\noindent{\textbf{Game $\game{#1}$:}}}
\newcommand{\conclude}{\medskip\noindent{}}
\newcommand{\myskip}{-0.16\baselineskip}

% Make vectors bold
\renewcommand{\vec}[1]{\ensuremath{\mathbf{#1}}}

% General mathematics
\newcommand{\range}[2] {[#1 \, .. \, #2]}
\newcommand{\SD}{\Delta}
\newcommand{\smallset}[1] {\{#1\}}
\newcommand{\bigset}[1] {\left\{#1\right\}}
\newcommand{\GRP} {\mathbb{G}}
\newcommand{\HRP} {\mathbb{H}}
%\newcommand{\pair} {\hat{e}}
\newcommand{\brak}[1] {\left(#1\right)}
\newcommand{\sbrak}[1] {(#1)}
\newcommand{\alg}[1] {\pcalgostyle{#1}}
\newcommand{\image} {\operatorname{im}}
\newcommand{\myland} {\,\land\,}
\newcommand{\mylor} {\,\lor\,}
\newcommand{\w}{\omega}
\newcommand{\const}{\pcpolynomialstyle{const}}
\newcommand{\p}[1]{\pcpolynomialstyle{#1}}
\newcommand{\ptx}{\p{t_{X^0}}}
\newcommand{\ptxsim}{\p{\tilde{t}_{X^0}}}
\newcommand{\ev}[1]{\widetilde{\pcpolynomialstyle{#1}}}
\newcommand{\maxconst}{\pcvarstyle{max}}
\newcommand{\noofc}{\numberofconstrains}
\newcommand{\noofw}{\pcvarstyle{m}}
\newcommand{\dconst}{\pcvarstyle{d}}
\newcommand{\multconstr}{\pcvarstyle{n}}
\newcommand{\linconstr}{\pcvarstyle{Q}}
\newcommand{\expected}[1]{\mathbb{E}\left[#1\right]}
\newcommand{\infrac}[2]{#1 / #2}
\newcommand{\eps}{\varepsilon}
\DeclareMathOperator{\SPAN}{Span}
\newcommand{\maxdeg}{\pcvarstyle{max}}
%polynomials
\newcommand{\pf}{\p{pf}}
\newcommand{\pF}{\p{F}}
\newcommand{\pa}{\p{a}}
\newcommand{\pb}{\p{b}}
\newcommand{\pc}{\p{c}}
\newcommand{\pz}{\p{z}}
\newcommand{\pt}{\p{t}}
\newcommand{\pR}{\p{R}}
\newcommand{\pr}{\p{r}}
\newcommand{\ptlo}{\p{t_{lo}}}
\newcommand{\ptmid}{\p{t_{mid}}}
\newcommand{\pthi}{\p{t_{hi}}}
\newcommand{\pS}{\p{S}}
\newcommand{\pT}{\p{T}}
\newcommand{\pta}{\tilde{\p{a}}}
\newcommand{\ptb}{\tilde{\p{b}}}
\newcommand{\ptc}{\tilde{\p{c}}}
\newcommand{\ptz}{\tilde{\p{z}}}
\newcommand{\ptZH}{\tilde{\p{Z_H}}}
\newcommand{\vX}{\vec{X}}
\newcommand{\vB}{\vec{B}}
\newcommand{\va}{\vec{a}}
\newcommand{\vb}{\vec{b}}
\newcommand{\vc}{\vec{c}}
\newcommand{\tzkproof}{\pi}

%reduction
\newcommand{\tb}{\tilde{b}} 
\newcommand{\tw}{\tilde{w}} 

% bilinear maps

\newcommand{\bmap}[2] {\left[#1\right]_{#2}}
\newcommand{\gone}[1] {\bmap{#1}{1}}
\newcommand{\gtwo}[1] {\bmap{#1}{2}}
\newcommand{\gi} {\iota}
\newcommand{\gtar}[1] {\bmap{#1}{T}}
\newcommand{\grpgi}[1] {\bmap{#1}{\gi}}
\newcommand{\gnone}[1]{\left[#1\right]}

\newcommand{\msg}[1]{\mathtt{#1}}

% zero knowledge
\newcommand{\oracleo}{\mathsf{O}}
\newcommand{\oraclesrs}{\oracleo_{srs}}
\newcommand{\oraclec}{\oracleo_\cdv}
% for srs updatability
\newcommand{\intent}{\msg{intent}}
\newcommand{\update}{\msg{update}}
\newcommand{\setup}{\msg{setup}}
\newcommand{\final}{\msg{final}}
\newcommand{\verifysrs}{\verify_{srs}}
\newcommand{\srsupdate}{\pcalgostyle{Update}}
\newcommand{\crs}{\pcvarstyle{crs}}
\newcommand{\srs}{\pcvarstyle{srs}}
%\newcommand{\td}{\pcvarstyle{td}}
\newcommand{\ip}[2]{\left\langle #1, #2\right\rangle}
\newcommand{\zkproof}{\pi}
\newcommand{\proofsystem}{\pcschemestyle{\Psi}}
\newcommand{\ps}{\proofsystem}
\newcommand{\psfs}{\proofsystem_\fs}
% \newcommand{\ps}{\proofsystem}
\newcommand{\nuppt}{\pcmachinemodelstyle{NUPPT}}
\newcommand{\ro}{\mathcal{H}}
\newcommand{\rof}[2]{\mathbf{\Omega}_{#1, #2}}
\newcommand{\trans}{\pcvarstyle{trans}}
\newcommand{\tr}{\pcvarstyle{tr}}
\newcommand{\instsize}{\pcvarstyle{n}}
\newcommand{\KG} {\mathsf{K}}
\newcommand{\kcrs} {\KG_{\crs}}
\newcommand{\fs}{\pcalgostyle{FS}}
\newcommand{\sigmaprot}{\pcalgostyle{\Sigma}}
\newcommand{\se}{\pcvarstyle{se}}
\newcommand{\snd}{\pcvarstyle{snd}}
\newcommand{\zk}{\pcvarstyle{zk}}


%rewinding---tree of transcripts
\newcommand{\pcboolstyle}[1]{\mathtt{#1}}
\newcommand{\treebuild}{\pcalgostyle{TreeBuild}}
\newcommand{\tree}{\pcvarstyle{T}}
\newcommand{\counter}{\pcvarstyle{counter}}


%PLONK related
\newcommand{\pcschemestyle}[1]{\mathsf{#1}}
\newcommand{\plonkprot}{\pcschemestyle{P}}
\newcommand{\plonkprotfs}{\pcschemestyle{P}_\fs}
\newcommand{\sonicprot}{\pcschemestyle{S}}
\newcommand{\sonicprotfs}{\pcschemestyle{S}_\fs}
\newcommand{\marlinprot}{\pcschemestyle{M}}
\newcommand{\marlinprotfs}{\pcschemestyle{M}_\fs}
\newcommand{\selector}[1]{\pcvarstyle{q_{#1}}}
\newcommand{\selmulti}{\selector{M}}
\newcommand{\selleft}{\selector{L}}
\newcommand{\selright}{\selector{R}}
\newcommand{\seloutput}{\selector{O}}
\newcommand{\selconst}{\selector{C}}
\newcommand{\chz}{\mathfrak{z}}
\newcommand{\ochz}{{\omega \mathfrak{z}}}
\newcommand{\reduction}{\rdv}
\newcommand{\tdv}{\pcadvstyle{T}}
\newcommand{\ch}{\pcvarstyle{ch}}

\newcommand{\game}[1]{\pcalgostyle{G}_{#1}}

\newcommand{\lag}{\p{L}}
\newcommand{\pubinppoly}{\p{PI}}

\newcommand{\ZERO}{\p{Z}}
%\newcommand{\pcsetstyle}[1]{\mathrm{#1}}
\newcommand{\HHH}{\mathsf{H}}
\newcommand{\KKK}{\mathsf{K}}

% general complexity theory
% \newcommand{\RND}[1]{\pcalgostyle{RND}(#1)}
\newcommand{\RND}[1]{\pcvarstyle{R}(#1)}
\newcommand{\RELGEN}{\mathcal{R}}
\newcommand{\REL}{\mathsf{REL}}
\newcommand{\LANG}{\mathcal{L}}
\newcommand{\inp}{\mathbbm{x}}
\newcommand{\impinp}{\mathbbm{y}}
\newcommand{\wit}{\mathbbm{w}}
\newcommand{\class}[1]{\mathfrak{#1}}
\newcommand{\ig}{\pcalgostyle{IG}}
\newcommand{\accProb}{\event{acc}}
\newcommand{\waccProb}{\event{\widetilde{acc}}}
\newcommand{\frkProb}{\event{frk}}
\newcommand{\extProb}{\event{ext}}
\newcommand{\ssndProb}{\event{ssnd}}
\newcommand{\FS}{\pcalgostyle{FS}} % Fiat-Shamir transform
\renewcommand{\aux}{\pcvarstyle{aux}} %auxiliary input
% \newcommand{\dlog}{\pcvarstyle{dlog}}
\newcommand{\vereq}{\p{ve}}

%Commitment schemes
\newcommand{\COM}{\pcschemestyle{C}}
\newcommand{\PCOM}{\pcschemestyle{PC}}
\newcommand{\PCOMp}{\pcschemestyle{PC}_{\plonkprot}}
\newcommand{\PCOMs}{\pcschemestyle{PC}_{\sonicprot}}
\renewcommand{\com}{\pcalgostyle{Com}}
\newcommand{\op}{\pcalgostyle{Op}}
\renewcommand{\open}{\op}

\newcommand{\committer}{\pcalgostyle{C}}
\newcommand{\receiver}{\pcalgostyle{R}}

%Plonk and Sonic
\newcommand{\ur}[1]{{#1\text{-}\mathsf{ur}}}

\newcommand{\plonk}{\ensuremath{\textnormal{\textsf{Plonk}}}}
\newcommand{\marlin}{{\ensuremath{\textnormal{\textsf{Marlin}}}}}
\newcommand{\sonic}{{\ensuremath{\textnormal{\textsf{Sonic}}}}}
\newcommand{\groth}{\ensuremath{\textsc{Groth16}}}
\newcommand{\plonkmod}{\ensuremath{\plonk^\star}}
\newcommand{\plonkint}{\ensuremath{\plonk^\star}}
\newcommand{\polyprot}{\pcalgostyle{poly}}
\newcommand{\plonkintpoly}{\plonkint_\polyprot}
% \newcommand{\sonic}{\textsc{Sonic}}
\newcommand{\maxdegree}{\pcvarstyle{N}}

\newcommand{\dlog}{\pcvarstyle{dlog}}
\newcommand{\ldlog}{\pcvarstyle{ldlog}}


%reductions
\newcommand{\extss}{\ext_\sfss}
\newcommand{\extse}{\ext_\se}
\newcommand{\extt}{\ext_{\pcvarstyle{tree}}}
\newcommand{\compass}{\mathtt{C}}
\newcommand{\ks}{\pcvarstyle{ks}}
\newcommand{\sfss}{\pcvarstyle{ss}}
\newcommand{\rdvks}{{\rdv_\ks}}
\newcommand{\rdvs}{{\rdv_\pcvarstyle{s}}}
\newcommand{\rdvdlog}{{\rdv_\dlog}}
\newcommand{\rdvldlog}{{\rdv_\ldlog}}
\newcommand{\rdvse}{{\rdv_\se}}
\newcommand{\rdvss}{{\rdv_\sfss}}
\newcommand{\rdvur}{\rdv_\pcvarstyle{ur}}

\newcommand{\env}{\pcadvstyle{E}}
\newcommand{\zdv}{\pcadvstyle{Z}}

\newcommand{\advse}{\adv}
\newcommand{\advss}{{\adv_\sfss}}
\newcommand{\bdvfs}{\bdv_{\FS}}

\newcommand{\epsss}{\eps_\pcvarstyle{f}}
\newcommand{\epsur}{\eps_\pcvarstyle{ur}}
\newcommand{\epsh}{\eps_{\pcvarstyle{hid}}}
\newcommand{\epsk}{\eps_{\pcvarstyle{k}}}
\newcommand{\epsbind}{\eps_\pcvarstyle{bind}}
\newcommand{\epsbinding}{\eps_\pcvarstyle{bind}}
\newcommand{\epsid}{\eps_{\pcvarstyle{id}}}
\newcommand{\epsop}{\eps_\pcvarstyle{op}}
\newcommand{\epss}{\eps_\pcvarstyle{s}}
\newcommand{\epsfor}{\eps_\pcvarstyle{f}}
\newcommand{\epsbatch}{\eps_\pcvarstyle{btch}}
\newcommand{\epsdlog}{\eps_\dlog}
\newcommand{\epsldlog}{\eps_\ldlog}
\newcommand{\epsuber}{\eps_{\pcvarstyle{uber}}}
\newcommand{\epszk}{\eps_{\pcvarstyle{zk}}}
\newcommand{\epsro}{\eps_{\ro}}

%selection polynomials
\newcommand{\vql}{\vec{q_{L}}}
\newcommand{\vqr}{\vec{q_{R}}}
\newcommand{\vqm}{\vec{q_{M}}}
\newcommand{\vqo}{\vec{q_{O}}}
\newcommand{\vx}{\vec{x}}
\newcommand{\vqc}{\vec{q_{C}}}

%errors
\newcommand{\err}{Err}
\newcommand{\errur}{\err_{ur}}
\newcommand{\errss}{\err_\sfss}
\newcommand{\errfrk}{\err_\frkProb}


%forking
\newcommand{\forking}{\pcalgostyle{F}}
\newcommand{\genforking}{\pcalgostyle{GF}}

%moving proofs and instances
\newcommand{\MoveInstanceForward}{\pcalgostyle{MoveInstanceForward}}
\newcommand{\MoveInstanceBackward}{\pcalgostyle{MoveInstanceBackward}}
\newcommand{\MoveProofForward}{\pcalgostyle{MoveProofForward}}
\newcommand{\MoveProofBackward}{\pcalgostyle{MoveProofBackward}}

%% Fields and other math
\newcommand{\HH}{\mathbb{H}}
\newcommand{\BB}{\{0,1\}}
\newcommand{\G}{\mathbb{G}}
\newcommand{\GT}{\mathbb{G}_T}
%Luca's commands and macros
\newcommand{\K}{\mathbb{K}}
\newcommand{\Q}{\mathbb{Q}}
\newcommand{\B}{\mathbb{B}}
\newcommand{\F}{\mathbb{F}}

%% hash, epsilons, protcools
\newcommand{\HASH}{\pcalgostyle{Hash}}
\newcommand{\vareps}{\varepsilon}
\newcommand{\varepscommit}{\vareps_{\text{commit}}}
\newcommand{\varepsquery}{\vareps_{\text{query}}}
\newcommand{\prot}{\pcalgostyle{\Pi}}
\newcommand{\protfs}{\pcalgostyle{\Pi}_{\FS}}
\newcommand{\proverfs}{\prover_{\FS}}
\newcommand{\verifierfs}{\verifier_{\FS}}

%colors
\definecolor{darkmagenta}{rgb}{0.5,0,0.5}
\definecolor{lightmagenta}{rgb}{1,0.85,1}
\definecolor{lightmagenta}{rgb}{0.9,0.9,0.9}
\definecolor{darkred}{rgb}{0.7,0,0}
\definecolor{blueish}{rgb}{0.1,0.1,0.5}
\definecolor{pinkish}{rgb}{0.9,0.8,0.8}
\definecolor{darkgreen}{rgb}{0,0.6,0}
\definecolor{lightgreen}{rgb}{0.85,1,0.85}
\definecolor{skyblue}{rgb}{0.3,0.9,0.99}
\definecolor{nicepurple}{rgb}{0.40784, 0.05490, 0.29412}

%\newcommand{\pcschemestyle}[1]{\mathbf{#1}}

\usepackage[most]{tcolorbox}

%comments
\DeclareRobustCommand{\markulf}[2]  {{\color{darkmagenta}\hl{\scriptsize\textsf{Markulf #1:} #2}}}
\DeclareRobustCommand{\michals}[2]  {{\color{blueish}{\scriptsize\textsf{Michal #1:}} #2}}
\newcommand{\chaya}[1]{\textcolor{magenta}{[{\footnotesize {\bf Chaya:} { {#1}}}]}}
\newcommand{\hamid}[1]{\textcolor{blue}{[{\footnotesize {\bf Hamid:} { {#1}}}]}}
\newcommand{\task}[2]{\todo[author=\textbf{Task},inline]{({\textit{#1}}) #2}}
% \newcommand{\task}[2] {\xcommenti{Task}{#1}{#2}}
% \DeclareRobustCommand{\task}[2]  {{\color{black}\sethlcolor{yellow}\hl{\textsf{TASK #1:} #2}}}
\DeclareRobustCommand{\changedm}[1] {{\color{violet} #1}}



% English Abbreviations
\usepackage{xspace}
\newcommand{\ie}{\text{i.e.}\xspace}
\newcommand{\etal}{\text{et al.}\xspace}
\newcommand{\naive}{\text{na\"ive}\xspace}
\newcommand{\wrt}{\text{w.r.t.}\xspace}
\newcommand{\whp}{\text{w.h.p.}\xspace}
\newcommand{\eg}{\text{e.g.}\xspace}
\newcommand{\cf}{\text{cf.}\xspace}
\newcommand{\visavis}{\text{vis-\`a-vis}\xspace}
\newcommand{\aka}{\text{a.k.a.}\xspace}
\newcommand{\ala}{\text{\`{a} la}\xspace}
\newcommand{\iid}[0]{\text{i.i.d.}\xspace}




\newcommand{\plonkytwoconstraint}{(\mathcal{P}_\inp, \mathcal{Q}_\inp, \sigma_\inp, \p{PI}=\emptyset, \p{r}_\inp, \p{r}', \p{\ell}_\inp, \p{n}_\inp,\p{t}, \FF, \mathbb{K})}

\newcommand{\plonkytwoconstraintnox}{(\mathcal{P}, \mathcal{Q}, \sigma, \p{PI}=\emptyset, \p{r}, \p{r}', \p{\ell}, \p{n},\p{t}, \FF, \mathbb{K})}


\newcommand{\partf}[2]{\bar{f}_{#1}(\omega^{#2})}
\newcommand{\partfx}[1]{\bar{f}_{#1}(X)}
\newcommand{\partfxk}[1]{\bar{f}_{k{#1}}(X)}
\newcommand{\partfxkfull}[2]{\bar{f}_{{#2}{#1}}(X)}
\newcommand{\partfxkfullx}[2]{\bar{f}_{{#2}{#1}}(x)}
\newcommand{\partffull}[3]{\bar{f}_{#1}(\tau_{1{#3}}, \omega^{#2})}
\newcommand{\partfxfull}[2]{\bar{f}_{#1}(\tau_{1{#2}},X)}


\newcommand{\partg}[2]{\bar{g}_{#1}(\omega^{#2})}
\newcommand{\partgx}[1]{\bar{g}_{#1}(X)}
\newcommand{\partgxk}[1]{\bar{g}_{k{#1}}(X)}
\newcommand{\partgxkfull}[2]{\bar{g}_{{#2}{#1}}(X)}
\newcommand{\partgxkfullx}[2]{\bar{g}_{{#2}{#1}}(x)}
\newcommand{\partgfull}[3]{\bar{g}_{#1}(\tau_{1{#3}}, \omega^{#2})}
\newcommand{\partgxfull}[2]{\bar{g}_{#1}(\tau_{1{#2}},X)}

\newcommand{\roundonemessage}{(\p{a}_i(X), \p{b}_i(X))_{i\in [\p{r}]}}
\newcommand{\roundonecha}{(\beta_i, \gamma_i)_{i\in [\p{t}]}}
\newcommand{\roundtwomessage}{(\p{z}_k(X), \pi_{ki}(X))_{k\in [\p{t}], i\in [\p{s}-1]}}
\newcommand{\roundtwocha}{(\alpha_i)_{i\in [\p{t}]}}
\newcommand{\roundonepartialtr}{ ((\p{a}_i(X),)_{i\in [\p{r}]}, (\beta_i, \gamma_i)_{i\in [\p{t}]})}
\newcommand{\roundonepartialtronebetagamma}[1]{ ((\p{a}_i(X))_{i\in [\p{r}]}, (\beta_{#1}, \gamma_{#1}))}
\newcommand{\roundtwopartialtr}{ ((\p{a}_i(X))_{i\in [\p{r}]}, \roundonecha, \roundtwomessage, \roundtwocha)_{i\in [\p{t}]})}
\newcommand{\roundtwopartialtronealpha}[1]{ ((\p{a}_i(X))_{i\in [\p{r}]}, (\beta_i, \gamma_i)_{i\in [\p{t}]}, (\p{z}_k(X), \pi_{ki}(X))_{k\in [\p{t}], i\in [\p{s}-1]}, \alpha_{#1})}
\newcommand{\polyroundtwofull}[1]{\p{d}(\inp, \tau_{2{#1}}, X)}
\newcommand{\polyroundtwofulleval}[2]{\p{d}(\inp, \tau_{2{#1}}, {#2})}

\newcommand{\pred}[1]{\mathsf{#1}}
\newcommand{\RS}{\pred{RS}}
\newcommand{\IPRS}{\pred{IPRS}}
\newcommand{\bbF}{\FF}
\renewcommand{\defeq}[0]{\ensuremath{{\;\vcentcolon=\;}}\xspace}

\newcommand{\idx}{\mathbbm{i}}
\newcommand{\permrelation}{\REL_{\mathbf{RPerm}}}
\newcommand{\permindex}{(\p{r},  \sigma)}

\newcommand{\turboplonkrelation}{{\REL_{\mathbf{RTurboPlonk}}}}
\newcommand{\plonkyrelation}{{\REL_{\mathbf{ROPlonky}}}}
\newcommand{\plonkytworelation}{{\REL_{\mathbf{RPlonky2}}}}

\newcommand{\oracle}[1]{\llbracket
{#1}\rrbracket
}
\newcommand{\incompleteoracle}[1]{\{\{#1\}\}}

\newcommand{\LT}{\p{LT}}
\newcommand{\turboplonkidx}{( \mathcal{P}, \mathcal{Q}, \sigma, \p{PI}, \p{r}, \p{\ell})}
\newcommand{\turboplonkconstraint}{\turboplonkidx}
\newcommand{\plonkyindex}{( \mathcal{P}, \mathcal{Q}, H, \sigma, \p{PI}, \p{r},\p{r}', \ell, \p{t})}
\newcommand{\plonkyindexexpanded}{(\FF, \mathbb{K}, \omega, \mathcal{P}, \mathcal{Q}, \sigma, \p{PI}, \p{r},\p{r}', \p{\ell},\p{t},\hash)}

\newcommand{\parttr}[1]{\mathcal{T}_{#1}}
\newcommand{\idxmaxf}{\p{idx_{max}}(F)}
\newcommand{\idxmaxg}{\p{idx_{max}}(G)}
\newcommand{\idxminf}{\p{idx_{min}}(F)}
\newcommand{\idxming}{\p{idx_{min}}(G)}
\newcommand{\roundonechanoidx}{(\beta, \gamma)}
\newcommand{\llex}{\leq_{\p{llex}}}
\newcommand{\pdiag}{\p{pdiag}}


\newcommand{\Ind}{\p{Ind}}


\newcommand{\plonky}{\p{Plonky}}
\newcommand{\plonkyindexer}{\Ind_{\plonky}}
\newcommand{\plonkyprover}{\prover_{\plonky}}
\newcommand{\plonkyverifier}{\verifier_{\plonky}}

\newcommand{\plonkyldtoracle}{\p{OPlonky}\xspace}
\newcommand{\plonkyldtindexer}{\Ind_{\plonkyldtoracle}}
\newcommand{\plonkyldtprover}{\prover_{\plonkyldtoracle}}
\newcommand{\plonkyldtverifier}{\verifier_{\plonkyldtoracle}}

\newcommand{\plonkyfs}[1]{\prot_{\p{Plonky2}, \FS}({#1})}
\newcommand{\plonkyproverfs}[1]{\prover_{\p{Plonky2}, \FS}({#1})}
\newcommand{\plonkyverifierfs}[1]{\verifier_{\p{Plonky2}, \FS}({#1})}


\newcommand{\plonkysmall}{\prot_{\p{rounds123}}}
\newcommand{\plonkysmallprover}{\prover_{\p{rounds123}}}
\newcommand{\plonkysmallverifier}{\verifier_{\p{rounds123}}}



\newcommand{\quotient}[3]{\p{quotient}({#1},{#2},{#3})}

\newcommand{\KK}{\mathbb{K}}

% \newcommand{\ro}{\mathsf{H}}
\newcommand{\zinc}{\mathsf{Zinc}}

\newcommand{\batchedfri}{\p{batchFRI}}
\newcommand{\ldtrelation}{\mathbf{PerfLDT}}
\newcommand{\ppp}{\p{gp}}
\newcommand{\pppc}{\mathsf{pp}_{\mathsf{Com}}}
\newcommand{\idxacc}{\idx_{\mathsf{acc}}}
\newcommand{\inpacc}{\inp_{\mathsf{acc}}}
\newcommand{\witacc}{\wit_{\mathsf{acc}}}

\newcommand{\permarg}{\prot_{\permrelation}}
\newcommand{\permargprover}{\prover_{\permrelation}}
\newcommand{\permargverifier}{\verifier_{\permrelation}}
\newcommand{\permargindexer}{\Ind_{\permrelation}}

\newcommand{\doomedldt}{\cD_{0}}
\newcommand{\doomedldtprox}{\cD_{\delta}}
\newcommand{\doomedldtks}{\cD_{0,\p{ks}}}
\newcommand{\doomedldtproxks}{\cD_{\delta,  \p{ks}}}

\newcommand{\union}{\ensuremath{\cup}\xspace}
\newcommand{\intersect}{\ensuremath{\cap}\xspace}
\newcommand{\pow}[1]{\ensuremath{^{(#1)}}\xspace}
\newcommand{\fin}{\pred{f}}
\newcommand{\tQ}{\ensuremath{\widetilde{Q}}\xspace}



\newcommand{\polytranscripts}[1]{\text{PolyTr}(#1)}
\newcommand{\mapstranscript}[1]{\text{Words}(#1)}
\newcommand{\extmapstranscript}[1]{\text{ExtWords}(#1)}


\newcommand{\epsfrirbr}{\ensuremath{\eps_{\p{rbr}}^{\p{FRI}}}\xspace}
\newcommand{\epsbatchfrirbr}{\ensuremath{\eps_{\p{rbr}}^{\p{bFRI}}}\xspace}
\newcommand{\epsfrirbrk}{\ensuremath{\eps_{\p{rbr-k}}^{\p{FRI}}}\xspace}
\newcommand{\epsbatchfrirbrk}{\ensuremath{\eps_{\p{rbr-k}}^{\p{bFRI}}}\xspace}
\newcommand{\epsfrifs}{\ensuremath{\eps_{\p{fs}}^{\p{FRI}}}\xspace}
\newcommand{\epsbatchfrifs}{\ensuremath{\eps_{\p{fs}}^{\p{bFRI}}}\xspace}
\newcommand{\epsfrifsq}{\ensuremath{\eps_{\p{fs-q}}^{\p{FRI}}}\xspace}
\newcommand{\epsbatchfrifsq}{\ensuremath{\eps_{\p{fs-q}}^{\p{bFRI}}}\xspace}

\newcommand{\?}{\ensuremath{\stackrel{?}{=}}\xspace}
\newcommand{\oplonky}{\ensuremath{\plonkyldtoracle}\xspace}


\newcommand{\oracleagreementdelta}{{\p{OCoAgg}(\delta)}}
\newcommand{\oracleagreementzero}{{\p{OCoAgg}(0)}}
\newcommand{\ldtrelationprox}[1]{{\mathbf{CoAgg({#1})}}}
\newcommand{\ldtrelationproxplain}{{\mathbf{CoAgg}}}
\newcommand{\protcoagg}[1]{\prot^{\p{OCoAgg}({#1})}}
\newcommand{\protcoaggactualrel}[1]{{\p{CoAgg}({#1})}}

\newcommand{\idxcorrelatedagreement}{(\FF, D, d, \delta, r)}


  \newcommand{\protcompiled}{ \prot_{\p{compiled}}  } \newcommand{\indexcompiled}{\Ind_{\p{compiled}}}
  \newcommand{\provercompiled}{\prover_{\p{compiled}}}  \newcommand{\verifiercompiled}{\verifier_{\p{compiled}}}

  \newcommand{\plonkyhIOP}{\p{Plonky2\-hIOP}}


    \newcommand{\ca}{\p{CA}}
    \newcommand{\idxinp}{\idx,\inp}


\newcommand{\rank}{\rho}
\newcommand{\Free}{\mathbb{FR}_\rank}
\newcommand{\Freenum}[1]{F_{#1}}
\newcommand{\Id}{\p{Id} }
\newcommand{\mle}[1]{{\widetilde{#1}}}
\newcommand{\hypercube}[1]{\{0,1\}^{#1}}
\newcommand{\FCCSunrolled}{(\rho, I=(f_1, \ldots, f_k), m,n,\ell,t,q,d), (M_0,\ldots, M_{t-1}, c_0,\ldots, c_{q-1}))}

%%ethSTARK
\newcommand{\traceevaldomain}{\pcvarstyle{D}}


\newcommand{\vect}[1]{\mathbf{#1}}
\newcommand{\bigtable}{\vect{t}}
\newcommand{\sizetable}{N}
\newcommand{\smalltable}{\vect{a}}
\newcommand{\smalltabletwo}{\vect{b}}
\newcommand{\smalltablethree}{\vect{c}}
\newcommand{\sizesmalltable}{m}
%\newcommand{\setvec}[1]{\left\{ #1 \right\}}

%%lookup tuples
\newcommand{\sizes}{(\sizetable, \sizesmalltable)}

\newcommand{\lookuptuple}{(\sizes, \bigtable; \smalltable)}

\newcommand{\indxedlookuptuple}{(\sizes, \bigtable; \smalltable, \tau)}

\newcommand{\indxedlookuptupleone}{(\sizes, \bigtable; \smalltable_1, \tau_1)}

\newcommand{\indxedlookuptupletwo}{(\sizes, \bigtable; \smalltable_2, \tau_2)}

\newcommand{\sizesnew}{(\sizetable+\sizesmalltable, \sizesmalltable)}

\newcommand{\indxedlookuptupleprime}{(\sizesnew, \bigtable'; \smalltable, \tau')}

\newcommand{\commatlooktuple}[1]{((N,m), \overline{M_{#1}}, \overline{\smalltable_{#1}}, \overline{\bigtable} ;  M_{#1}, \smalltable_{#1}, \bigtable)}

%%relations

\newcommand{\protsetequality}{\prot_{\p{set-eq}}}
\newcommand{\relsetequality}{\REL_{\p{set-equal}}}


    %\newcommand{\todo}[1]{}
    %\newcommand{\albert}[1]{}
    %\newcommand{\hendrik}[1]{}
    %\newcommand{\katy}[1]{}
    %\newcommand{\luca}[1]{}




%\newcommand{\ilia}[1]{\textcolor{nicepurple}{#1}}
%\newcommand{\nacho}[1]{\textcolor{orange}{#1}}
%\newcommand{\matteo}[1]{\textcolor{red}{#1}}
%\newcommand{\nick}[1]{\textcolor{violet}{#1}}



\newcommand{\protperm}{\prot_{\p{perm}}}
\newcommand{\relperm}{\REL_{\p{perm}}}
\newcommand{\tildepol}[1]{\widetilde{#1}}
\newcommand{\ind}{\mathbb{I}}

\newcommand{\Rccs}{\REL_{\p{cR1CS}}}
\newcommand{\Rcs}{\REL_{\p{R1CS}}}
\newcommand{\Rcrcs}{\REL_{\p{acc}}}
\newcommand{\Rcsa}{\Rcrcs(\textbf{r})}
\newcommand{\Rms}{\REL_{\p{ms}}}
\newcommand{\Rsps}{\REL_{\p{sps}}}

%%variables
\newcommand{\bb}{\vec{b}}
\newcommand{\yy}{\vec{y}}
\newcommand{\xx}{\vec{x}}
\newcommand{\XX}{\vec{X}}
\newcommand{\EE}{\mathbf{E}}
\newcommand{\zz}{\vec{z}}
\renewcommand{\ZZ}{\mathbb{Z}}

\renewcommand{\ss}{\vec{s}}
\renewcommand{\rr}{\vec{r}}
\newcommand{\YY}{\mathbf{Y}}
\newcommand{\WW}{\mathbf{W}}

%\newcommand{\lEE}{\lin(\mathbf{E})}

\newcommand{\TT}{\mathbf{T}}
\newcommand{\multilinmle}[1]{\widetilde{#1}}
\newcommand{\multilin}{\mathsf{ml}}

\newcommand{\bWW}{{\widebar{\WW}}}
\newcommand{\bEE}{{\widebar{\EE}}}
\newcommand{\bTT}{\widebar{\TT}}



\newcommand{\MM}{\mathbf{M}}
\newcommand{\bMM}{\widebar{\mathbf{M}}}
%%lookup relations
%%normal lookup
\newcommand{\rellookup}{\REL_{\mathsf{Look}}}
\newcommand{\relsoslookup}{\REL_{\mathsf{SOSLook}}}
\newcommand{\relindexedlookup}{\REL_{\p{IdxLook}}}
%%algebraic lookup
\newcommand{\relalookup}{\REL_{\p{AlLook}}}
\newcommand{\relalookupeq}{\REL_{\p{EqAlLook}}}
\newcommand{\relaindexedlookup}{\REL_{\p{AlIdxLook}}}
\newcommand{\relccl}{\REL_{\p{ConstAlCmLook}}}
%%committed lookup relations
\newcommand{\relMlookupcom}{\REL_{\p{CmMAlLook}}}
%%shift relations
\newcommand{\relshift}{\REL_{\mathsf{Shift}}}
\newcommand{\relshifttwo}{\REL_{\mathsf{Shift2}}}
\newcommand{\relshiftbatch}{\REL_{\mathsf{ShiftBatch}}}

\newcommand{\cm}{\p{cm}}

\newcommand{\meq}{\multilinmle{\p{eq}}}
\newcommand{\miota}{\multilinmle{\iota}}
\newcommand{\mtau}{\multilinmle{\tau}}
\newcommand{\Racc}{\mathcal{R}_{\mathrm{acc}}}
\newcommand{\Relax}{\mathcal{R}_{\mathrm{rR1CS}}}
\newcommand{\Reval}{\mathcal{R}_{\mathrm{eval}}}
\newcommand{\Rsmol}{\mathcal{R}_{\mathrm{decR1CS}}^b}
\newcommand{\Fp}{\mathbb{F}_p}
\newcommand{\nin}{n_{in}}
\newcommand{\R}{\mathcal{R}}
\newcommand{\Rp}{\mathcal{R}_p}
\newcommand{\C}{\mathcal{C}}
\newcommand{\assgn}{\xleftarrow{}}
\newcommand{\flatten}[1]{#1^\flat}
\newcommand{\deflatten}[1]{#1^\sharp}

\renewcommand{\norm}[1]{\left\lVert#1\right\rVert_{\infty}}
\newcommand{\norminf}[1]{\norm{#1}_\infty}
\newcommand{\normop}[1]{\norm{#1}_{\mathrm{op}}}
\newcommand{\Ajt}{\mathcal{A}jt}
\newcommand{\NTT}{\p{NTT}}
\newcommand{\RotSum}{\p{RotSum}}
\newcommand{\bsmol}{(1,\:b,\ldots,\:b^{k-1})}
\newcommand{\vvec}{\p{vec}}
\newcommand{\draw}{\xleftarrow{\$}}
\newcommand{\Csmall}{\mathcal{C}_{\mathrm{small}}}
\newcommand{\stark}{2^{251}+17\cdot2^{192}+1}
\newcommand{\circlep}{2^{31}-1}
\newcommand{\babybear}{15 \cdot 2^{27} + 1}
\newcommand{\goldilocks}{2^{64} - 2^{32}  + 1}

\newcommand{\ronecs}{\p{R1CS}}
\newcommand{\cronecs}{\p{cR1CS}}
\newcommand{\crronecs}{\p{crR1CS}}
\newcommand{\red}{\p{red}}
\newcommand{\iop}{\p{IOP}}
\newcommand{\I}{\mathcal{I}}

\newcommand{\Mod}[1]{\ (\mathrm{mod}\ #1)}
\newcommand{\ord}{\mathrm{ord}}
\newcommand{\N}{\mathbb{N}}
\newcommand{\Z}{\mathbb{Z}}
\newcommand{\inter}{\langle \Prover^*, \Verifier \rangle}
\newcommand{\hinter}{\langle \Prover, \Verifier \rangle}
%To have widebar command without adding the whole package
\DeclareFontFamily{U}{mathx}{\hyphenchar\font45}
\DeclareFontShape{U}{mathx}{m}{n}{
      <5> <6> <7> <8> <9> <10>
      <10.95> <12> <14.4> <17.28> <20.74> <24.88>
      mathx10
      }{}
\DeclareSymbolFont{mathx}{U}{mathx}{m}{n}
\DeclareFontSubstitution{U}{mathx}{m}{n}
\DeclareMathAccent{\widebar}{0}{mathx}{"73}



\newcommand{\blue}[1]{\textcolor{blue}{#1}}
\newcommand{\lin}{\mathcal{L}}

\newcommand{\Requal}{\REL_{\mathsf{equal}}}


\newcommand{\Prover}{\prover}
\newcommand{\Verifier}{\verifier}


\newcommand{\relsumcheck}{\REL_{\p{sc}}}
\newcommand{\releval}{\REL_{\p{eval}}}


\newcommand{\Rccsnew}{\REL_{\mathsf{cR1CS}}}
\newcommand{\Rrccs}{\REL_{\mathsf{rcR1CS}}}
\newcommand{\Rrcs}{\REL_{\mathsf{rR1CS}}}

\newcommand{\gen}{\mathsf{Gen}}
\newcommand{\pparams}{\mathsf{gp}}


    \newcommand{\eabort}{\mathcal{E}_{\p{abort}}}
        \newcommand{\evalidwit}{\mathcal{E}_{\p{valid-wit}}}
    \newcommand{\pro}{\p{Pr}}
    \newcommand{\expectation}{\mathbb{E}}
    \newcommand{\cE}{\mathcal{E}}
            \newcommand{\cElines}{\cE_{\mathsf{lines}}}

                \newcommand{\eone}{\mathcal{E}_{\p{root}}^{\EE^{(1)}}}
        \newcommand{\etwo}{\mathcal{E}_{\p{good-run}}}


            \newcommand{\UU}{\mathbf{U}}
    \newcommand{\VV}{\mathbf{V}}
		\newcommand{\cEzero}{\cE_{\p{zero}}}
		\newcommand{\cEnonzero}{\cE_{\p{nonzero}}}
		\newcommand{\cEbinding}{\cE_{\p{binding}}}

  \newcommand{\cEbadlines}{\mathcal{E}_{\p{bad-lines}}}


  \renewcommand{\gp}{\mathsf{gp}}
\newcommand{\vp}{\mathsf{vp}}
\newcommand{\idxer}{\mathcal{I}}

\renewcommand{\oracle}[1]{[[{#1}]]}
\newcommand{\loc}[1]{\ZZ_{(#1)}}
\newcommand{\vv}{\mathbf{v}}
\newcommand{\cC}{\mathcal{C}}
\newcommand{\cR}{\mathcal{R}}

\newcommand{\cL}{\mathcal{L}}
\newcommand{\lift}{\pcvarstyle{lift}}
\newcommand{\bounds}{\ppp}


\newcommand{\cQ}{\mathcal{Q}}
\newcommand{\cF}{\mathcal{F}}
\newcommand{\soundness}{\mathsf{sound}}
\newcommand{\boundedsoundness}{\mathsf{bound-sound}}

\newcommand{\completeness}{\mathsf{comp}}
\newcommand{\completenesserror}{\epsilon}
%\renewcommand{\multilin}{\mathsf{multilin}}

\newcommand{\nummorph}{\eta}
\newcommand{\cA}{\mathcal{A}}
\newcommand{\valid}{\mathsf{wf}}
\newcommand{\Enc}{\mathsf{Enc}}

\newcommand{\codelength}{m}
\newcommand{\codedim}{\ell}
\newcommand{\uu}{\mathbf{u}}
\newcommand{\T}{\mathsf{T}}
\newcommand{\bdsamplingset}{S}
\newcommand{\cP}{\mathcal{P}}
\renewcommand{\qq}{\mathbf{q}}
\newcommand{\local}[1]{\ZZ_{(#1)}}
\newcommand{\bounded}{\mathsf{bounded}}
\newcommand{\rate}{\rho}
\newcommand{\test}{\mathsf{test}}
\newcommand{\distance}{\mathsf{dist}}
\newcommand{\dd}{\mathbf{d}}
\newcommand{\accept}{\mathsf{accept}}
\newcommand{\bdpoly}{\mathsf{bdpoly}}
\renewcommand{\aa}{\mathbf{a}}
\newcommand{\Runtime}{\mathsf{Runtime}}
\newcommand{\Encbinary}{\mathsf{Enc}_{\mathsf{binary}}}
\newcommand{\wt}{\mathsf{wt}}
\newcommand{\supp}{\mathsf{supp}}

\newcommand{\Ber}{\mathsf{Ber}}
\newcommand{\BP}{\mathsf{BP}}
\newcommand{\rk}{\mathsf{rk}}
\newcommand{\sgn}{\mathsf{sgn}}

    \newcommand{\PC}{\mathsf{PC}}
    \newcommand{\LA}{\mathsf{LA}}
    \newcommand{\prove}{\mathsf{Prove}}


    \newcommand{\advA}{\mathcal{A}}
    \newcommand{\Commit}{\mathsf{Commit}}
    \newcommand{\Open}{\mathsf{Open}}
    \newcommand{\Eval}{\mathsf{Eval}}
    \newcommand{\binding}{\mathsf{binding}}



\DeclareMathOperator{\Frac}{Frac}
\newcommand{\lcm}{\text{lcm}}

\newcommand{\Ext}{\mathsf{Ext}}

\newcommand{\record}{\mathsf{Record}}
%\newcommand{\ks}{\mathsf{ks}}
\newcommand{\ww}{\mathbf{w}}
%  \renewcommand{\Span}{{\mathsf{span}}}
\newcommand{\li}{\mathsf{l.i.}}
\newcommand{\ld}{\mathsf{l.d.}}
\newcommand{\boundgenerator}{\lVert M_\cC \rVert_\infty}
\newcommand{\CCS}{\mathsf{CCS}}
\newcommand{\RICS}{\mathsf{R1CS}}
\newcommand{\ZZZ}{\mathbf{Z}}
\newcommand{\WWW}{\mathbf{W}}
\newcommand{\www}{\mathbf{w}}

   % Define macros for the procedures
    \newcommand{\Setup}{\mathsf{Setup}}
    \newcommand{\Com}{\mathsf{Com}}
    \newcommand{\ProveEvalMod}{\mathsf{ProveEvalMod}}
    \newcommand{\VfyEvalMod}{\mathsf{VfyEvalMod}}
    
    % Define macros for common symbols
    \newcommand{\opn}{\mathsf{opn}}
    \newcommand{\cmt}{\mathsf{cm}}

\newcommand{\BDtestsamplingprime}{{q_0}}
\newcommand{\codelin}{{\cC_{\lambda}}}
\newcommand{\boundq}{{B_{\mathsf{pt}}}}
\newcommand{\boundp}{\boundq}

\newcommand{\boundv}{{B_{\vv}}}

\newcommand{\primeset}{{\cP}}
\newcommand{\proj}{{\mathsf{proj}}}
\newcommand{\total}{\mathsf{total}}

\newcommand{\extraction}{\mathsf{extraction}}
\renewcommand{\time}{\mathsf{time}}

\newcommand{\codelinplain}{\cC}
\newcommand{\good}{\mathsf{good}}
\newcommand{\len}{\mathsf{len}}
\newcommand{\witnessbound}{{2\cdot\boundv + (6\codedim+2)\cdot \log(\BDtestsamplingprime\cdot \codedim)}}
%\newcommand{\witnessbound}{{\boundv + \log(\BDtestsamplingprime^3\cdot \codedim)}}
\newcommand{\witnessboundencoded}{\log(\norm{M_\codelin}\cdot \codedim) +\witnessbound }

\newcommand{\pos}{\mathsf{pos}}
\renewcommand{\LANG}{\mathsf{LANG}}

% Bit-vector helpers (used in Appendix/extended_preliminaries.tex)
\newcommand{\Val}{\mathsf{Val}}
\newcommand{\Next}{\mathsf{Next}}
\newcommand{\hi}{\mathsf{hi}}
\newcommand{\bx}{\mathbf{x}}
\newcommand{\by}{\mathbf{y}}
\newcommand{\bu}{\mathbf{u}}

\newcommand{\cM}{\mathcal{M}}
\newcommand{\relevalppp}{(\mu, \boundv, \boundp, \delta, \BDtestsamplingprime)}

\newcommand{\words}{\mathsf{Words}}

\newcommand{\thetaexpression}{\frac{2\cdot \codedim^2\cdot \left( \boundv +\log(\codedim)  + 2 \cdot \log(\BDtestsamplingprime)\right) + \boundv +\mu\cdot \boundp  + 3\mu+2}{\lambda\cdot |\primeset|}}

\newcommand{\thetaexpressionnotreduced}{\frac{ 2\cdot\codedim^2\cdot \left( \log(\norm{\words(\hat{\uu}})) + 2 \cdot \log(\BDtestsamplingprime)\right) + \boundv +\mu\cdot \boundp  + 3\mu+2}{\lambda\cdot |\primeset|}}

\newcommand{\lo}{{\mathsf{lo}}}

\newcommand{\tot}{{\mathsf{to}}}
\newcommand{\inpy}{\mathbbm{y}}

\newcommand{\ex}{{\mathsf{ex}}}
%\newcommand{\size}{{\mathsf{size}}}
\newcommand{\mul}{{\mathsf{mul}}}
\newcommand{\samp}{{\mathsf{sample}}}
\newcommand{\jea}{{\mathsf{JEA}}}
\newcommand{\ns}{{\mathsf{ns}}}

\newcommand{\cD}{\mathcal{D}}
\renewcommand{\tt}{\mathbf{t}}
\newcommand{\Zinc}{\mathsf{Zinc}}
\newcommand{\Zip}{\mathsf{Zip}}
\newcommand{\zipp}{\textsf{Zip+}\xspace}



\newcommand{\PIOP}{\mathsf{PIOP}}
\newcommand{\IOP}{\mathsf{IOP}}

\newcommand{\inproof}{\mathsf{IP}}
\newcommand{\indexer}{\mathsf{I}}
\newcommand{\vparams}{\mathsf{vp}}
\newcommand{\gparams}{\mathsf{gp}}

\newcommand{\bA}{\mathbf{A}}
\newcommand{\bB}{\mathbf{B}}
\newcommand{\bC}{\mathbf{C}}

\renewcommand{\pp}{\mathbf{p}}

\newcommand{\hint}{\mathsf{hint}}

\newcommand{\weirdrationals}{\QQ}%$[\hat{\uu}]}

\newcommand{\Irr}{\mathsf{Big}}

 \newcommand{\symbolbigbound}{\mathsf{BoundEnc}}

 \newcommand{\compiled}{\mathsf{compiled}}
 \newcommand{\succinct}{\mathsf{succ}}
 \newcommand{\vcom}{\mathsf{VC}}
 \newcommand{\JEA}{\mathsf{JEA}}
  \newcommand{\weight}{\mathsf{weight}}
  \newcommand{\dom}{\mathcal{D}}
  \newcommand{\eq}{\mathsf{eq}}
  \newcommand{\gparamssigma}{\mathsf{gp}^{\Sigma}}

%  \newcommand{\band}{\ \mathsf{AND}\ }
\newcommand{\xor}{\ \mathsf{XOR}\ }
%\newcommand{\eq}{\mathsf{eq}}


\newcommand{\lookupset}{S}
\newcommand{\lookupsets}{\mathcal{S}}
\renewcommand{\set}{\mathsf{set}}
\renewcommand{\bin}{\widetilde{\mathsf{\mathbf{bin}}}}

\renewcommand{\int}{\mathsf{int}}
\newcommand{\lagrange}{\mathsf{Lag}}

% Notation symbols (so \Word_n and \BitPoly_n work)
\newcommand{\Word}{\mathsf{Word}}
\newcommand{\BitPoly}{\mathsf{BitPoly}}

% Optional “set-valued” versions (so you can write \Wordset{n}, \BitPolyset{n})
\newcommand{\Wordset}[1]{\Range{0}{2^{#1}-1}}
\newcommand{\BitPolyset}[1]{\bits^{<#1}[X]}
% or: \newcommand{\BitPolyset}[1]{\set{\sum_{i=1}^{#1} b_i X^{i-1}\mid b_i\in\bits}}
\newcommand{\rot}{\mathsf{ROT}}


\newcommand{\megaQ}{\mathfrak{Q}}
\renewcommand{\ff}{\mathbf{f}}
\newcommand{\cS}{\mathcal{S}}
\newcommand{\gpunrolled}{(k, m, n, \mu, B, d, t)}

\newcommand{\ucs}{\REL_{\mathsf{UCS}}}
\newcommand{\witlen}{m}
\newcommand{\inplen}{k}
\newcommand{\degreebound}{d}
\newcommand{\bitbound}{B}
\newcommand{\polbounds}{(\degreebound,\bitbound)}
\newcommand{\constraints}{\mathcal{C}}
\newcommand{\primetuple}{\mathbf{q}}
\newcommand{\constraintpol}{Q}
\newcommand{\numconstraints}{|\constraints|}
\newcommand{\degbounds}{\mathbf{d}}
%\newcommand{\numrows}{n}
\newcommand{\rowvar}{K}
\newcommand{\rowindex}{k}
\newcommand{\primeindex}{i}
\newcommand{\constraintindex}{t}
\newcommand{\numprimes}{|\primetuple|}
\newcommand{\colindex}{j}
\newcommand{\inpvar}{\YY}
\newcommand{\numcols}{c}

\newcommand{\relmodN}{\REL_{\mathsf{MOD}}}
\newcommand{\relNF}{\REL_{\mathsf{Ring}}}
\newcommand{\evalconstraints}{\text{\emph{Eval}}}
\newcommand{\mainQpolyring}{(\QQ[X]_{\leq B, d})}
\newcommand{\mainpolyring}[3]{#1^{<{#3}, {#2}}[X]}
\newcommand{\mainQpolyringmultilin}{(\QQ[X]_{\leq B, d})^{\multilin}}

\newcommand{\albert}[1]{\textcolor{MidnightBlue}{Albert: {#1}}}
\newcommand{\albertimportant}[1]{\textcolor{red}{Albert: {#1}}}
\newcommand{\ilia}[1]{\textcolor{nicepurple}{#1}}

\newcommand{\UUU}{\mathbf{U}}
\newcommand{\SSS}{\mathbf{S}}
\newcommand{\MMM}{\mathbf{M}}
\newcommand{\PPP}{\mathbf{P}}
\newcommand{\QQQ}{\mathbf{Q}}
\newcommand{\RRR}{\mathbf{R}}
\newcommand{\AAA}{\mathbf{A}}
\providecommand{\band}{\ \mathsf{AND}\ }

\newcommand{\ROTR}{\mathsf{ROTR}}
\newcommand{\ROTL}{\mathsf{ROTL}}
\newcommand{\SHR}{\mathsf{SHR}}
\newcommand{\Ch}{\mathsf{Ch}}
\newcommand{\Maj}{\mathsf{Maj}}


\newcommand{\tab}{\mathbf{t}}
\newcommand{\tablelen}{\tau}
\newcommand{\relronecs}{\REL_{\mathsf{R1CS}}}
\newcommand{\uair}{\REL_{\mathsf{UAIR+}}}
\newcommand{\uairproj}{\REL_{\mathsf{UAIR+}}^{\mathsf{proj}}}

\newcommand{\setintegers}{S}
\newcommand{\numintegers}{2^s}
\newcommand{\ones}{\mathbf{1}}
\newcommand{\getsr}{\gets}

\newcommand{\rem}{\mathsf{rem}}


\newcommand{\rd}{\mathsf{rd}}
\newcommand{\po}{\mathsf{p}}
\newcommand{\rdp}{\rd_{\po}}
\newcommand{\wh}{\widehat}
\newcommand{\ucsalg}{\REL_{\mathsf{UCS-ALG}}}
\newcommand{\ucsalgproj}{\REL_{\mathsf{UCS-ALG-PROJ}}}
\newcommand{\Kk}{\mathbb{K}}
\newcommand{\Qb}{\mathsf{Q}}
\newcommand{\qs}{\mathsf{q}}
\newcommand{\X}{\mathbf{X}}
\newcommand{\sr}{\mathsf{sr}}
\newcommand{\trsr}{\mathsf{tr}_{\mathsf{sr}}}


\newcommand{\powertuples}{\bm{\ell}}
\newcommand{\evalpoints}{\mathbf{a}}

\renewcommand{\proj}{\mathsf{proj}}
\newcommand{\id}{\mathsf{id}}
\newcommand{\Bad}{\mathsf{Bad}}

\newcommand{\numchunks}{\mathsf{c}}
\renewcommand{\sizetable}{N}
\renewcommand{\bigtable}{\mathbf{t}}
\newcommand{\den}{\mathsf{den}}%
\renewcommand{\num}{\mathsf{num}}%

\newcommand{\protuair}{\Pi_{\mathsf{UAIR}^+}}

% Commands moved from PIOP/
\newcommand{\ucsm}{\mathsf{UCS}}
\newcommand{\epsdef}{\eps_{\mathsf{def}}}
\newcommand{\epswf}{\varepsilon_{\mathsf{wf}}}
\newcommand{\epsproj}{\varepsilon_{\mathsf{proj}}}
\newcommand{\acc}{\mathcal{E}_{\mathsf{acc}}}
\newcommand{\rat}{\mathsf{rat}}
\newcommand{\albertforkai}[1]{\textcolor{purple}{Albert for Kai: {#1}}}

\newcommand{\bfg}{\mathbf{g}}
\newcommand{\bfh}{\mathbf{h}}
\newcommand{\ftwoz}{{$\mathbb{B}\text{it}\ZZ$}}

\newcommand{\comfieldexp}{\nu}
\newcommand{\pack}{\mathsf{Pack}_\comfieldexp}
\newcommand{\comfield}{\FF_{2^\comfieldexp}}
%\newcommand{\evalfield}{\FF_{q}}
\newcommand{\evalfield}{\KK}
\newcommand{\bff}{\mathbf{f}}
\renewcommand{\oracle}{\mathbbm{y}}
\renewcommand{\Enc}{\mathsf{Enc}}
\newcommand{\picomfield}{\pi_{2}}
\newcommand{\pievalfield}{\pi_q}
\newcommand{\basegenerator}{g}
\newcommand{\wordsize}{{32}}
\newcommand{\degevalfield}{\mathsf{deg}(\KK)}


% MACROS

\renewcommand{\enspace}{\ensuremath{\thinspace}}
% operators
\newcommand{\deq}{\ensuremath{\mathrel{\coloneq}}}


% objects
\newcommand{\FExt}{\ensuremath{\mathbb{H}}}
\newcommand{\CX}{\ensuremath{\mathbb{C}}}
\newcommand{\E}{\ensuremath{\mathbb{E}}}
\newcommand{\pmset}{\ensuremath{\{-1,1\}}}

\newcommand{\Bits}{\{0,1\}}

% malicious
\newcommand{\Malicious}[1]{\tilde{#1}}

% language and relation

\newcommand{\ShortIndex}{\iota}
\newcommand{\InstanceExplicit}{\mathbbm{x}}
\newcommand{\InstanceImplicit}{\mathbbm{y}}
\newcommand{\Witness}{\mathbbm{w}}
\newcommand{\Proof}{\pi}
\newcommand{\GetLanguage}[1]{\Language(#1)}

\newcommand{\Indexer}{\mathbf{I}}
\newcommand{\Decider}{\mathbf{D}}
\newcommand{\Extractor}{\mathbf{E}}



\newcommand{\Tiny}[1]{{\scriptscriptstyle #1}}
\newcommand{\Length}{\mathsf{l}}
\newcommand{\IndexerLength}{\mathsf{l}_{\Tiny{\IndexerProof}}}
\newcommand{\ProofLength}{\Length}
\newcommand{\Queries}{\mathsf{q}}
\newcommand{\Randomness}{\mathsf{r}}
\newcommand{\Round}{\mathsf{k}}
\newcommand{\Knowledge}{\kappa}
\newcommand{\Alphabet}{\Sigma}

\newcommand{\SecParam}{\lambda}
\newcommand{\RepetitionParameter}{t}
\newcommand{\Interaction}[2]{\langle #1, #2 \rangle}

\newcommand{\HammingDistanceOp}{\Delta}
\newcommand{\Distance}{\delta}
\newcommand{\HammingBall}[1]{B_{#1}}
\newcommand{\HammingBallVolume}[1]{\mathsf{Vol}_{#1}}
\newcommand{\ReedSolomon}{\mathsf{RS}}
\newcommand{\InterleavedReedSolomon}{\mathsf{IRS}}
\newcommand{\Rate}{\rho}
\newcommand{\MessageSize}{\ell}
\newcommand{\CodeSize}{n}
\newcommand{\EvaluationDomain}{\mathcal{L}}
\newcommand{\AmbientSpace}{\Z^\CodeSize}
\newcommand{\ErasureCorrector}{\Extractor_{\Code}}
\newcommand{\Code}{\mathcal{C}}
\newcommand{\CodeInverse}{\Code^{\Tiny{-1}}}
\newcommand{\LinearCodeDefinition}{\Code\colon \F^{\MessageSize} \to \F^{\CodeSize}}
\newcommand{\CodeEncoding}{\mathsf{enc}_{\Code}}
\newcommand{\CodeGenerator}{G_{\Code}}
\newcommand{\CodeSuccinctEvaluation}{\mathsf{ev}_{\Code}}
\newcommand{\CodePerfectDecoding}{\mathsf{dec}_{\Code}}
\newcommand{\CodeErasureDecoding}{\mathsf{ecor}_{\Code}}
\newcommand{\CodeDefinition}{\Code\colon \Alphabet^{\MessageSize} \to \Alphabet^{\CodeSize}}
\newcommand{\CodeInterleaved}[2]{{#1^{\equiv_{#2}}}}
\newcommand{\MinimumDistance}[1]{\Distance_{\Tiny{\mathsf{min}}}(#1)}
\newcommand{\Codeword}{u}
\newcommand{\CodewordAlternative}{v}
\newcommand{\ListSize}[2]{\left|\List\left(#1, #2\right)\right|}
\newcommand{\ListSizeCode}{L_{\Code}(\Distance)}


\newcommand{\StateWitness}{\mathsf{w}}
\newcommand{\KnowledgeState}{\State}
\newcommand{\KnowledgeStateWitness}{\StateWitness}
\newcommand{\Func}{f}

\newcommand{\RandomnessCombine}{\gamma}
\newcommand{\EQPoly}{\mathsf{eq}}
\newcommand{\SumcheckPoly}{\hat{h}}
\newcommand{\Target}{\mu}
\newcommand{\SumcheckTarget}{\Target}

\newcommand{\FSize}{|\F|}

\newcommand{\ErrorMutualCA}{\varepsilon_{\Tiny{\mathsf{mca}}}}
\newcommand{\ErrorCA}{\varepsilon_{\Tiny{\mathsf{ca}}}}

%Succinct Evaluation Stuff
\newcommand{\SuccinctEvaluation}{\mathsf{V}}
\newcommand{\SuccinctEvaluationState}{\mathsf{st}}
\newcommand{\SuccinctEvaluationDomain}{\mathcal{S}_{\SuccinctEvaluation}}
\newcommand{\SuccinctEvaluationVec}{\mathsf{ev}}
\newcommand{\FixSE}[2]{\mathsf{fix}_{#2}(#1)}
\newcommand{\OutputPairing}{\beta_{\Tiny{\mathsf{out}}}}
\newcommand{\OutputSuccinctEvaluation}{\SuccinctEvaluation_{\Tiny{\mathsf{out}}}}
\newcommand{\OutputSuccinctEvaluationState}{\SuccinctEvaluationState_{\Tiny{\mathsf{out}}}}
\newcommand{\BatchedTarget}{\Target_{\Tiny{\mathsf{batch}}}}

% Protocol params
\newcommand{\FoldingParameter}{k}
\newcommand{\ZeroEvader}{\mathsf{Z}}
\newcommand{\ZeroEvaderDomain}{\mathcal{S}_{\Tiny{\mathsf{ze}}}}
\newcommand{\ZeroEvaderLength}{\ell_{\Tiny{\mathsf{ze}}}}
\newcommand{\ZeroEvaderSeed}{s_{\Tiny{\mathsf{ze}}}}
\newcommand{\ErrorZeroEvader}{\varepsilon_{\Tiny{\mathsf{ze}}}}
\newcommand{\ErrorProjection}{\varepsilon_{\Tiny{\mathsf{proj}}}}
\newcommand{\ProjectedRing}{R_{\Phi}}
\newcommand{\ErrorSumcheck}[1]{\varepsilon_{\Tiny{\mathsf{sc}}}(#1)}
\newcommand{\BatchingZeroEvader}{\mathsf{B}}
\newcommand{\BatchingZeroEvaderDomain}{\mathcal{S}_{\Tiny{\mathsf{batch}}}}
\newcommand{\BatchingZeroEvaderSeed}{s_{\Tiny{\mathsf{batch}}}}
\newcommand{\ErrorBatching}{\varepsilon_{\Tiny{\mathsf{batch}}}}


%\newcommand{\PIMSAT}{\Relation_{\mathsf{PIMSAT}}}
\newcommand{\moduledegree}{d}
\renewcommand{\ideal}{\mathcal{I}}
\newcommand{\numrows}{n}

\newcommand{\ringbitbound}[1]{\ring_{\leq{#1}}}
\newcommand{\ringdegbound}[1]{\ring^{<{#1}}}
\newcommand{\ringtwobounds}[2]{\ring_{\leq {#1}}^{<{#2}}}

\newcommand{\hamming}{\Delta}

\newcommand{\LinBitsDual}{\Relation_{\mathsf{Lin-BitsDual}}}
%\newcommand{\LinModFieldBin}{\Relation_{\mathsf{Lin-ModFieldsBin}}}
\newcommand{\LinModFieldBin}{\Relation_{\mathsf{Lin-BitsDual}}}

\renewcommand{\red}{\mathsf{red}} 
\newcommand{\iopp}{\mathsf{iopp}} 

\newcommand{\ksred}{{\kappa_{\red}}}
\newcommand{\ksiopp}{{\kappa_{\iopp}}}
\newcommand{\LIN}{\Relation_{\mathsf{Lin}}}
\newcommand{\Mmatrix}{A}

\newcommand{\ringtwo}{R_{\mathsf{e}}}

\newcommand{\moduleone}{{M_\ring}}
\newcommand{\moduletwo}{{\ringtwo}}
\newcommand{\ringmorphism}{\pi}
\newcommand{\modulemorphism}{\psi}
\newcommand{\moduleonebound}{{M_{\ringbitbound{B}}}}

\newcommand{\picanon}{\pi_{\mathsf{canon}}}

\newcommand{\gkrpol}{p}
\renewcommand{\cc}{\mathbf{c}}
\newcommand{\RELCM}{\Relation_{\mathsf{CM}}}


\newcommand{\UPESAT}{\Relation_{\mathsf{UPESAT}}}
\newcommand{\PIMSAT}{\UPESAT} % old name of UPESAT, kept as an alias

\newcommand{\FFext}{\FF_{\mathsf{ext}}}




\newcommand{\PESAT}{\Relation_{\mathsf{PESAT}}}

\newcommand{\virtual}{\mathsf{virtual}}

%
\newcommand{\ronecsideal}{\Relation_{\mathsf{R1CS-UPESAT}}}
\newcommand{\iopmlechallenge}{S}

\newcommand{\fieldboundsizeexponent}{(|\comfield|-1)/\codedim_1}

\newcommand{\bitstext}{\mathsf{bits}}
\newcommand{\pirand}{\pi_{\mathsf{rand}}}


\newcommand{\Char}{\mathsf{char}}
\renewcommand{\dim}{\mathsf{dim}}



\title{\plaintitle}

\author[1]{}

\affil[1]{}


\date{\today}


%%%%%%%%%%%%%%%%%%%%%%%%%%%%%%%%%%%%%%%%%%%%%%%%%%%%%%%%%%%%%%%%%%%%%%%%%%%%%%%%
\begin{document}
\maketitle

\begin{abstract}
\albert{Should the name be $\FF 2 \FF$?}
\end{abstract}

\setcounter{tocdepth}{3}

\tableofcontents
%%%%%%%%%%%%%%%%%%%%%%%%%%%%%%%%%%%%%%%%%%%%%%%%%%%%%%%%%%%%%%%%%%%%%%%%%%%%%


\section{Introduction}

\subsection{Acknowledgements}
%World Foundation, Lev Soukhanov, Remco, Ulrich, Giacomo, Marcin

\section{Preliminaries}


Given a prime $q$ we let $\pi_q: \ZZ\to \FF_q$ be the ring homomorphism mapping an integer to its reduction modulo $q$.  We further let $\pi_0$ denote the identity map on $\ZZ$, so that $\pi_0(x) = x$ for all $x\in \ZZ$. We let $\FF_0 = \ZZ$.

Given a  finite field $\FF$ of characteristic $q$, let $a(X)\in \ZZ[X]$ be an irreducible polynomial over $\FF_q$ \albert{yes?}, so that $\FF = \ZZ[X]/(q,a(X))$ \albert{consider isomorphism, or assume all finite fields are represented like this}. Let $\psi_{\FF}: \ZZ[X] \to \FF$ be the  canonical homomorphism from $\ZZ[X]$ to $\FF$ obtained by reducing modulo the ideal $(q,a(X))$, i.e.\ by reducing coefficients modulo $q$ and then reducing the resulting polynomial modulo $a(X)$.




\paragraph{Interactive Oracle Proofs of Proximity (IOPP)}
Interleaved codes

\paragraph{Interactive Oracle Reductions (IOR)}


\section{\ftwoz}

%\newcommand{\comfield}{\FF_{\mathsf{c}}}
%\newcommand{\evalfield}{\FF_{\mathsf{e}}}

%\newcommand{\psicomfield}{\psi_{\mathsf{c}}}
%\newcommand{\psievalfield}{\psi_{\mathsf{e}}}


\subsection{Constrained codes over dual fields \albert{not sure dual is a good word: the two fields are just ``different'', while dual things tend to have a close relation between each other}}
Let $\codedim_0\geq 0$ be a dimension parameter, and $\codelength\geq 0$ a blocklength parameter. Let $\FF$ be a field (not necessarily a prime field), and let $\cC\subseteq \FF^{\codelength}$ be a linear code of dimension $\codedim_0$ and blocklength $\codelength$ over $\FF$.  Let $\inplen\geq0$ be a public input length parameter, and let $Q\in \FF[Y_1, \ldots, Y_{\inplen}, Z_1, \ldots, Z_{\codedim}]$  be a constraint polynomial, of arbitrary degree, 
with coefficients in $\FF$ and on $\inplen+\codedim$ variables. Further, let $\codedim\geq \codedim_0$ be a witness length paramter, and let $P:\FF^{\codedim} \to \FF^{\codedim_0}$ be an arbitrary map (looking ahead, for us $P$ will be a packing map that packs elements from $\FF_2$ into elements of a larger binary field.     Consider the following relation:
%
$$
\Relation^Q[\cC,\FF, P] \defeq \left\{ \left(\begin{aligned} &\inp=\yy\in \FF^\inplen, \\&\oracle = f \in \FF^\codelength, \\ &\wit= \bff\in \FF^\codedim\end{aligned}\right) \;\middle|\;  \begin{aligned} &\Enc_{\cC}(P(\bff)) = f\\ \wedge \  &Q(\yy, \bff) = 0  \end{aligned} \right\}
$$
When $Q$ and $\yy$ are such that $\yy=(\vv,\mu)\in \FF^{\codedim}\times \FF$ and $$Q(\yy,\ff)= \langle \vv, \ff\rangle = \sum_{i\in [\codedim]} \vv_i \cdot \ff_i,$$ we denote $\Relation^Q[\cC,\FF, P]$  by $\Relation^{\mathsf{lin}}[ \cC,\FF,P]$.

%When $Q$ and $\yy$ are such that $\yy=(\rr,\mu)\in \FF^{\log\codedim}\times \FF$ and $$Q(\yy,\ff)= \mle{\bff}(\rr)-\mu,$$ we denote $\Relation_{\cC, \FF}^{Q}$  by $\Relation_{\cC, \FF}^{\mathsf{mle}}$. 


Next let $\comfield$ be a finite field of characteristic $2$,  and let $\evalfield$ be  a finite field of prime order $q$ or $\QQ$. For intuition, $\comfield$ is the field over which we will commit our witness, and $\evalfield$ is the field over which we prove linear relations on the witness. We let
%
$$
\psicomfield : \ZZ \to \FF_2 \subseteq \comfield, \qquad \psievalfield: \ZZ \to \evalfield,
$$
be reduction mod $2$ and $q$ maps. We futher let 
$$
\psicomfield^{-1} : \FF_2 \to \{0,1\} \subseteq \ZZ, \qquad \psievalfield^{-1}: \evalfield \to \range{0}{q-1} \subseteq \ZZ,
$$
be maps such that $\psicomfield(\psicomfield^{-1}(x))=x$ and $\psievalfield(\psievalfield^{-1}(x))=x$ for all $x$. We remark that one could choose any other such ``inverse'' maps, e.g.\ for $\psievalfield^{-1}$ one may prefer to use centered representatives. 



Assume $|\comfield|=2^\comfieldexp$. Let
%
$$
\beta_1, \ldots, \beta_\comfieldexp \in \comfield
$$
%
be an $\FF_2$-base of $\comfield$ seen as an $\FF_2$-vector space of dimension $n$.

Given a vector $\uu\in \FF_2^\codelength$, with $2^\comfieldexp$ dividing $\codelength$, we define
%
$$
\pack(\uu)\in \comfield^{\codelength/\comfieldexp}
$$
as the vector whose $j$-th entry is
%
$$
\pack(\uu)_{j} = \uu_{(j-1)\cdot \comfieldexp +1 }\cdot \beta_1 + \ldots + \uu_{(j-1)\cdot \comfieldexp +\comfieldexp-1 }\cdot \beta_{\comfieldexp-1}+\uu_{j \cdot \comfieldexp} \cdot \beta_\comfieldexp \in \comfield.
$$
%
%
Let $\cC\subseteq \comfield^{\codelength}$ be a linear code of dimension $\codedim/2^{\comfieldexp}$ and blocklength $\codelength$ over $\comfield$. Define
$$
\Relation^{\mathsf{lin-dual}}[\cC,\comfield, \evalfield,\pack] \defeq \left\{ \left(\begin{aligned} &\inp=(\mu, \vv)\in \evalfield \times \evalfield^{\codedim}, \\&\oracle = f \in \comfield^\codelength, \\ &\wit= \bff\in \{0,1\}^\codedim \subseteq \ZZ^\codedim\end{aligned}\right) \;\middle|\;  \begin{aligned} &\Enc_{\cC}(\pack(\psicomfield(\bff))) = f\\ \wedge \  & \langle \psievalfield(\bff), \vv\rangle  = \mu  \end{aligned} \right\}
$$
%
\subsection{Tensor decomposition, interleaved codes, and products of constrained codes}
%
We fix a tensor decomposition of vectors $\vv\in \comfield^\codedim$ as $$\vv = \vv^{(1)} \otimes \vv^{(2)} = \left(\vv_i^{(1)}\cdot \vv_j^{(2)}\right)_{i\in [\codedim_1], j\in [\codedim_2]}$$ for some $\vv^{(1)}\in \evalfield^{\codedim_1}, \vv^{(2)}\in \evalfield^{\codedim_2}$, with $\codedim_1 \cdot \codedim_2=\codedim$. Note that the trivial tensor decomposition is admitted, i.e.\ we may effectively not decompose $\vv$ by choosing $\codedim_1=\codedim$, $\codedim_2=1$, $\vv^{(1)}=\vv$, and $\vv^{(2)} = (1)$.

We assume $\cC = \cC_1^{\equiv \codedim_2}$ is the $\codedim_2$-dimensional interleaving of a code $\cC_1\subseteq \comfield^{\codelength/\codedim_2}$ of dimension $\codedim_1/2^\comfieldexp$ and blocklength $\codelength/\codedim_2$.  We then index the entries of $\bff$ as $\bff=(\bff_{ij})_{i\in [\codedim_1], j\in[\codedim_2]}$. Note that
%
$$
\langle \psievalfield(\bff), \vv \rangle = \sum_{j\in [\codedim_2]}\left(\sum_{i\in [\codedim_1]} \psievalfield(\bff_{ij}) \cdot \vv_i^{(1)}\right) \cdot \vv_j^{(2)} = \sum_{j\in [\codedim_2]} \langle \psievalfield(\bff_{*j}) , \vv^{(1)} \rangle \cdot \vv_j^{(2)},
$$
where $\bff_{*j} = (\bff_{ij})_{i\in[\codedim_1]}$.








% We use $0$ and $1$ indistinctly to denote the zero and one integer, and zero and one field elements. 


Define the following bivariate polynomial:
$$
P(Y,Z) = (Y-1)\cdot Z +1 
$$
 Note
%
$
P(Y,0) = 1$ and $P(Y,1) = Y.
$
%
Consider the relation $\Relation_{\comfield, \pack}^{Q}$ with underlying code $\cC_1$ and 
%
$$Q(Y_1, \ldots, Y_{\codedim_1+1}, Z_1, \ldots, Z_{\codedim_1}) = \Pi_{i=1}^{\codedim_1} P(Y_i, Z_i) - Y_{\codedim_1+1}.$$ 
%
Precisely, 
%
$$
\Relation^{Q}[ \cC_1,\comfield, \pack] = \left\{ \left(\begin{aligned} &\inp=(\yy,z)\in \comfield^{\codedim_1}\times \comfield, \\&\oracle = f \in \comfield^{\codelength/\codedim_2}, \\ &\wit= \bff\in \comfield^{\codedim_1}\end{aligned}\right) \;\middle|\;  \begin{aligned} &\Enc_{\cC_1}(\pack(\bff)) = f\\ \wedge \  &\Pi_{i=1}^{\codedim_1} \left((\yy_i-1)\cdot \bff_i +1\right) = z  \end{aligned} \right\}
$$
%
We assume prover and verifier have access to an IOR $\Pi_{\mathsf{red}}$  that reduces 
%
$
\Relation^{Q}[\cC_1,\comfield,\pack]$ to $\Relation^{\mathsf{lin}}[ \cC_1,\comfield, \pack]$. 
%
%
We assume $\Pi_{\mathsf{red}}$ works so that the verifier never queries $\oracle$ \albert{Would be better to say upfront what this protocol is in our mental model, and then make abstract assumptions like this, which look arbitrary otherwise. Same applies below}.

Further, we consider the $\codedim_2$-fold parallel composition $\Pi_{\mathsf{red}}^{\codedim_2}$ of $\Pi_{\mathsf{red}}$,  and we assume that given $\codedim_2$ triples 
%
$$
\left((\yy^{(j)},\mu^{(j)}) , g^{(j)}, \bfh^{(j)}  \right)_{j\in[\codedim_2]} \in \left(\Relation^{Q}[ \cC_1,\comfield, \pack]\right)^{\codedim_2}
$$
the IOR $\Pi_{\mathsf{red}}^{\codedim_2}$ outputs 
%
$$
\left((\vv,\mu_j) \in \comfield^{ \codedim_1 } \times \comfield, g^{(j)} , \bfh^{(j)} \right)_{j\in[\codedim_2]} \in \left(\Relation^{\mathsf{lin}}[ \cC_1,\comfield, \pack]\right)^{\codedim_2},
$$
for some $\rr\in \comfield^{\log \codedim_1}$ and some $\mu_j\in \comfield$ for $j\in [\codedim_2]$.
In words, $\Pi_{\mathsf{red}}^{\codedim_2}$ outputs $\codedim_2$ linear claims for the input witnesses $\bfh^{(j)}$ behind the input oracles $g^{(j)}$, all \emph{on the same vector of coefficients  $\vv$}.   \albert{We could be more abstract with this IOR, but I find it difficult to tie it with the protocol in the next section then. Maybe better to write a remark} 

We describe one such $\Pi_{\mathsf{red}}^{\codedim_2}$ in \cref{?}, where $\vv$ is a vector of the form $\vv=(\mathsf{eq}(\bb,\rr))_{\bb\in \{0,1\}^{\log \codedim_1}}$, for some $\rr\in \comfield^{\log \codedim_1}$.


\subsection{The main {\ftwoz}  IOR}
We provide an IOPP for $\Relation^{\mathsf{lin-dual}}[\cC, \comfield, \evalfield, \pack]$ next.
\begin{construction}[IOPP for {$\Relation^{\mathsf{lin-dual}}[\cC, \comfield, \evalfield, \pack]$}]\label{c:core_iop}
The prover and verifier receive $\inp=(\mu, \vv)\in \evalfield \times \evalfield^{\codedim}$ and  $\oracle= f \in \comfield^\codelength$. The prover receives $\wit= \bff\in \bits^\codedim \subseteq \ZZ^\codedim$ such that $(\inp,\oracle, \wit) \in \Relation^{\mathsf{lin-dual}}[\cC, \comfield, \evalfield, \pack]$.  We assume the tensor split from the previous Section \ref{}, and adopt its notation. 
Let $\basegenerator\in \comfield$ be a nonzero element from $\comfield$ with sufficiently large multiplicative order, to be determined later.  


\begin{enumerate}[leftmargin=*]
\item The prover sends $\mu_1, \ldots, \mu_{\codedim_2}\in \ZZ$ such that $\langle \bff_{*j}, \psievalfield^{-1}(\vv^{(1)})\rangle = \mu_j$. 
\item The verifier checks that $\sum_{j\in [\codedim_1]} \psievalfield(\mu_j)\cdot \vv_j^{(2)}  = \mu$ over $\evalfield$.  \albert{We could have the prover commit to the $\mu_j$ and have this check proved by the PCS (a PCS that can prove linear constraints like WHIR). This pays off for very large $\codedim$.}
\item Prover and verifier define $\bm{\gamma}_i = \psievalfield^{-1}(\vv_i^{(1)}) \in \range{0}{q-1}\subseteq \ZZ$ for $i \in [\codedim_1]$, and set to prove the following claim:
%
$$
\prod_{i \in [\codedim_1]} g^{\bm{\gamma}_i\cdot \bff_{ij}}  = g^{\mu_j} \in \comfield, \qquad \text{for all } j\in [\codedim_2].
$$
Note that this claim is equivalent to
$$
\prod_{i \in [\codedim_1]} \left(\left(g^{\bm{\gamma}_i}-1\right)\cdot \bff_{ij} +1\right)= g^{\mu_j}  \quad  \in \comfield, \qquad \text{for all } j\in [\codedim_2].
$$
% \ulrich{%
% Suggestion: 
% Maybe it is more clear, if we write in the above two equalities explicitly a $\bmod\: q$, or, add $\in \mathbb F_q$?
% }%
% \albert{See now. It's actually in $\comfield$}
To prove the above claim, prover and verifier engage in the IOR $\Pi_{\mathsf{red}}^{\codedim_2}$   with input \albert{The reason I'm leaving this IOR abstract here is that I think there may be better ways to prove the above claim than a grand product GKR, so I want to write it modular} $$\left(\inp_j=((g^{\bm{\gamma}_i})_{i \in [\codedim_1]},\ g^{\mu_j}),\  \oracle_j= f_{*j}, \ \wit_j=\psicomfield(\bff_{*j})\right)_{j\in [\codedim_2]},$$
%
where in the honest case  $f_{*j} = \Enc_{\cC_1}(\pack(\psicomfield(\bff_{*j})))\in \comfield^{\codelength/\codedim_2}$.
%
 At the end of this IOR, they are left with the following claim to prove: 
%
\begin{equation}\label{e:iopp_mle_claims}
\left((\vv, \mu_j'), \ f_{*j}, \ \psicomfield(\bff_{*j})\right)_{j\in [\codedim_2]} \in \left(\Relation^{\mathsf{lin}}[\cC_1, \comfield,\pack]\right)^{\codedim_2}\end{equation} for some vector $\vv$ and field elements $\mu_j'$ output during the IOR $\Pi_{\mathsf{red}}^{\codedim_2}$. 
\item Prover and  verifier engage in a IOPP  to certify the claims \eqref{e:iopp_mle_claims}. 
\end{enumerate}
\end{construction}

\begin{theorem}\label{t:thm_core_IOPP}
    \albert{Completeness and RBR knowledge soundness of \cref{c:core_iop} }
\end{theorem}
\begin{proof}
It follows from \cref{t:thm_virtual_IOPP} by taking $\codedim_1'=\codedim_1$ and taking $M$ to be the identity matrix.
\end{proof}

We now describe a GKR-based IOR reducing $(\Relation_{\cC_1, \comfield}^{Q})^{\codedim_2}$ to  $(\Relation_{\cC_1, \comfield}^{\mathsf{eval}})^{\codedim_2}$ of the required form.

\begin{construction}[Grand product GKR]
\albert{Todo: it's a standard grand product GKR so far -- though there's lots of implementation-level optimizations it can undergo}
\end{construction}
\ulrich{%
Maybe to just write verbally, that we will use GKR for the grand product.
And defer the explicit protocol description to the appendix or any later section?
At this point it is maybe more important to clarify how useful the construction is because of packed commitment scheme over $\mathbb F_2$. 
% At Albert: should I omit such comments here in the source, and rather post them on slack?
}
\albert{Sounds good}


\section{Virtualization of $\FF_2$-linear combinations lifted to integers}


In this section we expand relation $\Relation^{\mathsf{lin-dual}}[\cC, \comfield, \evalfield, \pack]$ to impose, instead of constraints of the form $\Enc_\cC(\pack(\pi_2(\bff)))=f$ and $\langle \pi_q(\bff), \vv  \rangle = \mu$, constraints as $\Enc_\cC(\pack(\pi_2(\bff))) = f$  and $$\langle \pi_q(\pi_2^{-1}(M\cdot \pi_2(\bff))) \rangle = \mu$$ for a matrix $M\in \FF_{2}^{\codedim' \times \codedim}$. Precisely, given a new parameter $\codedim'$, we define 
%
$$
\Relation^{\mathsf{lin-virtual}}[\cC, \comfield, \evalfield, \pack] \defeq \left\{ \left(\begin{aligned} &\inp=(\mu, \vv, M)\in \evalfield \times \evalfield^{\codedim'} \times \FF_2^{\codedim' \times \codedim}, \\&\oracle = f \in \comfield^\codelength, \\ &\wit= \bff\in \bits^\codedim \subseteq \ZZ^\codedim\end{aligned}\right) \;\middle|\;  \begin{aligned} &\Enc_{\cC}(\pack(\psicomfield(\bff))) = f\\ \wedge \  & \langle \psievalfield(\pi_2^{-1}(M \cdot \pi_2(\bff)), \vv\rangle  = \mu  \end{aligned} \right\}
$$
%
The relation has the tag ``lin(ear)-virtual'' to reference the intuition that  $\Relation^{\mathsf{lin-virtual}}[\cC, \comfield, \evalfield, \pack] $ expresses a linear constraint on the vector $\bfh=\psievalfield(\pi_2^{-1}(M \cdot \pi_2(\bff))\in \evalfield^{\codedim'}$, which is not committed directly: only $\pi_2(\bff)$  is. Note that $\bfh$ is a linear expression over $\comfield$ of the elements in $\pi_2(\bff)$, lifted to the integers, and then projected to $\evalfield$.

\albert{Would be good to add more detail in this paragraph} The relation $\Relation_{\cC, \comfield, \evalfield}^{\mathsf{lin-virtual}}$  allows to, for example, express linear constraints over $\evalfield$ of the XOR of witness elements, while only committing (over $\comfield$) to the original elements. The same applies with other bit operations such as AND and Majority, by noting first that, given integer bits $x,y,z\in \{0,1\}$: \albert{Should reference some identity in preliminaries, or explain more here} 
%
\begin{equation*}
    \begin{aligned}
        \pi_q(x \band y) & = 2^{-1}\cdot \pi_q \left( x+y - \pi_2^{-1}(\pi_2(x) + \pi_2(y)) \right),\\
        \pi_q(\Maj(x,y,z)) &= 2^{-1}\cdot \pi_q \left( x+y+z - \pi_2^{-1}(\pi_2(x) + \pi_2(y) + \pi_2(z) \right),
    \end{aligned}
\end{equation*}
where $2^{-1}$ denotes the inverse of $2$ modulo $q$.  
%
Hence AND and Majority can be expressed over $\FF_q$ as  linear expressions of a vectors of the form $\pi_q(\pi_2^{-1}(M\cdot \pi_2(\bff)))$, with $\bff$ being a vector of integer bits containing only the original elements on which we apply the AND and Majority operations. 


\bigskip

We next provide an IOPP for $\Relation^{\mathsf{lin-virtual}}[\cC, \comfield, \evalfield, \pack]$ for specific shapes of $M$. Similarly as in \cref{?}, we assume a tensor structure $\codedim=\codedim_1 \cdot \codedim_2$ for $\vv$, and write $\bff = (\bff_{*j})_{j\in [\codedim_2]} \in \bits^{\codedim_1 \times \codedim_2}$.  We then let $\codedim_1'$ be such that $\codedim' = \codedim_1' \times \codedim_2$ and assume $M = \mathsf{Id} \otimes M_0$ where $M_0$ is a $\codedim_1' \times \codedim_1$ matrix with entries in $\FF_2$, and $\mathsf{Id}$ is the $\codedim_2\times \codedim_2$ matrix. In words, $M$ is the $\codedim'\times \codedim = \codedim_1'\cdot \codedim_2 \times \codedim_1\cdot \codedim_2$ diagonal block matrix with $M_0$ as diagonal blocks. Then
%
\begin{equation}\label{e:M_pi_f}
M \cdot \pi_2(\bff) = (M_0 \cdot \pi_2(\bff_{*j}))_{j\in [\codedim_2]} \in (\FF_2^{\codedim_1'})^{\codedim_2}.
\end{equation}
%
\begin{construction}[IOPP for $\Relation_{\cC, \comfield, \evalfield}^{\mathsf{lin-virtual}}$]\label{c:virtual_iop}
The prover and verifier receive $\inp=(\mu, \vv,M)\in \evalfield \times \evalfield^{\codedim'} \times \FF_2^{\codedim'\times\codedim}$ and  $\oracle= f \in \comfield^\codelength$. The prover receives $\wit= \bff\in \bits^\codedim \subseteq \ZZ^\codedim$ such that $(\inp,\oracle, \wit) \in \Relation_{\cC, \comfield, \evalfield}^{\mathsf{lin-virtual}}$.  We assume $M$ is of the form described above.

 Define  $\bfh = \pi_2^{-1}\left(M\cdot \pi_2(f)\right) \in \{0,1\}^{\codedim'} \subseteq \ZZ^{\codedim'}$. By \eqref{e:M_pi_f}, $$\bfh=(\bfh_{*j})_{j\in [\codedim_2]} = \left(\pi_2^{-1}\left( M_0\cdot \pi_2(\bff_{*j}) \right)\right)_{j\in [\codedim_2]},$$ with $ \bfh_{*j}\in \bits^{\codedim_1'}$. Let $f= \Enc_{\cC}(\pack(\pi_2(\bff))) \in \bits^{\codelength} \subseteq \ZZ^{\codelength}$.
 \begin{enumerate}[leftmargin=*]
     \item Prover and verifier execute \cref{c:core_iop}  with inputs 
$$
\inp' = (\mu, \vv),  \quad \oracle'=f,\quad  \wit' = \bfh=\pi_2^{-1}\left(M\cdot \pi_2(\bff)\right) \in \{0,1\}^{\codedim'} \subseteq \ZZ^{\codedim'}.
$$
 until Step 3 where they are meant to use the IOR $\Pi_{\mathsf{red}}$. Note that until that point, the verifier has not made any queries to $\oracle'$. So, in terms of completeness, it does not matter, yet, that $f$ is not the encoding of $\pi_2(\bfh)=M\cdot \pi_2(\bff)$.
 \item  Prover and verifier engage in $\Pi_{\mathsf{red}}^{\codedim_2}$ with inputs
 %
 $$\left(\inp_j=((g^{\bm{\gamma}_i})_{i \in [\codedim_1']},\ g^{\mu_j}),\  \oracle_j= f_{*j}, \ \wit_j=\psicomfield(\bfh_{*j})=M_0\cdot \pi_2(\bff_{*j})\right)_{j\in [\codedim_2]},$$
 where $f_{*j}=\Enc_{\cC_1}(\pack(\pi_2(\bff_{*j})))$. At the end of $\Pi_{\mathsf{red}}^{\codedim_2}$ they are left with claims
 %
%
\begin{equation}\label{e:virtual_IOPP_intermediate}
\left((\vv, \eta_j), \ \oracle_j, \ \wit_j\right)_{j\in [\codedim_2]} \in \left(\Relation^{\mathsf{lin}}[\cC_1, \comfield,\pack]\right)^{\codedim_2},
\end{equation}
%
for some $\vv\in \comfield^{ \codedim_1'}$ and $\eta_j \in \comfield$. Again, since we assume the verifier never queries $\oracle_j$ during $\Pi_{\mathsf{red}}$, completeness holds for honest provers, even if $\oracle_j$ is not the encoding of $\wit_j$.

Note that \eqref{e:virtual_IOPP_intermediate} holds if and only if, for all $j\in [\codedim_2]$,
%
\begin{align*}
\eta_j &=\sum_{i\in [\codedim_1']} \vv_i \cdot \pi_2(\bfh_{ij}) = \sum_{i\in [\codedim_1']} \vv_i \cdot \sum_{t\in [\codedim_1]} M_{0it}\cdot \pi_2(\bff_{tj})\\ =& \sum_{t\in [\codedim_1]} \left(\sum_{i\in[\codedim_1']}M_{0it}\cdot \vv_i\right)\cdot \pi_2(\bff_{tj}) = \langle \uu_j , \pi_2(\bff_{*j})\rangle,
\end{align*}
where $\uu \in \comfield^{\codedim_1}$ is the vector whose $t$-th entry is
%
$$
\uu_t = \langle \vv, M_{0*t}\rangle \in \comfield,
$$
and  $M_{0*t}\in \comfield^{\codedim_1'}$ denotes the $t$-th column of $M_0$, for $t\in [\codedim_1]$,
%

%
Note that, if $\Pi_{\mathsf{red}}^{\codedim_2}$ is as in \cref{?}, then $\vv= (\mathsf{eq}(\bb,\rr))_{\bb\in \{0,1\}^{\log \codedim_1'}}$ for some $\rr\in \comfield^{\log \codedim_1'}$, and then 
$
\uu_t = \mle{M_{0*t}}(\rr) \in \comfield. 
$
%

\item If the prover is honest, we have that
%
\begin{equation}\label{e:final_virt_equation}
\left((\uu, \eta_j),\ \oracle_j = \Enc_{\cC_1}(\pack(\pi_2(\bff_{*j}))), \ \pi_2(\bff_{*j})\right)_{j\in [\codedim_2]} \in \left(\Relation^{\mathsf{lin}}[\cC_1, \comfield,\pack]\right)^{\codedim_2}.
\end{equation}
%
Prover and verifier use an IOPP to prove \eqref{e:final_virt_equation}.

\albert{To keep in mind: ring-switch for arbitrary innter products could be expensive}
%
 \end{enumerate}
%
\end{construction}


\begin{example}
    Say we have a vector $\bff\in \bits^{\codedim} \subseteq \ZZ^\codedim$ and we wish to prove MLE evaluation claims of an extended vector $\bff' \in \bits^{2\cdot \codedim}$ where 
    %
    $$
    \bff_i' = \begin{cases}
    \bff_i  & \text{if } i \in [\codedim],\\
    \bff_{i -\codedim } \xor \bff_{i - \codedim +1} & \text{if } \codedim+1 \leq i \leq 2\cdot \codedim -1,\\
    \bff_{\codedim} & \text{if } i =2\cdot \codedim.
    \end{cases}
    $$
    % 
    In words, $\bff'$ is the result of concatenating $\bff$ and the vector obtained by ``xor-ing'' together every two consecutive elements in $\bff$ (and keeping the last entry as $\bff_{\codedim}$, since $\bff_{\codedim} \xor \bff_{\codedim+1}$ is not defined). This example will come up in our SHA-256 case-study.

    
    We have
    %
    $$
    \bff' = \pi_2^{-1}\left(M \cdot \pi_2(\bff)\right)   \in \bits^{2\codedim}
    $$
    where
    %
    $$
M = \begin{pmatrix}
        & & & &\\
       & & \mathsf{Id}_{\codedim\times \codedim} & & \\
               & & & &\\
               \hline
        1 & 1 &  0 & 0 & \ldots & 0\\
        0 & 1 & 1 & 0 & \ldots & 0 \\
        \vdots & & \ddots & & & \vdots\\
        0 &  \ldots &  0 & 0 & 1 & 1 \\
        0 & \ldots & 0& 0 & 0 & 1 \\
    \end{pmatrix} 
    $$
     if we do not wish to leverage any tensor decomposition structure, i.e.\ if we take $\codedim_1=\codedim$ and $\codedim_2=1$, then we can apply  \cref{c:virtual_iop} using this matrix $M$ as part of the input. At the end of \cref{c:virtual_iop} prover and verifier are left with the claim 
     $$
     \left((\uu, \eta),\ \Enc_{\cC}(\pack(\pi_2(\bff))), \ \pi_2(\bff)\right) \in \Relation^{\mathsf{lin}}[\cC_1, \comfield,\pack],
     $$
     for  $\uu = (\mle{M_{*t}}(\rr))_{t\in [\codedim]} \in \comfield^{\codedim}$, where $M_{*t}$ denotes the $t$-th column of $M$. I.e.\ prover and verifier must now certify that $\Enc_{\cC}(\pi_2(\bff))$ is (close to) the encoding of a vector $\pi_2(\bff)$ such that
     %
     $$
    \langle (\mle{M_{*t}}(\rr))_{t\in [\codedim]}, \pi_2(\bff) \rangle = \eta.
     $$
    %
    Using \cref{s:ring_switch}, \eqref{?} 


    
    Next we address the case 
     
\end{example} 

\begin{theorem}\label{t:thm_virtual_IOPP}
    \albert{Completeness and RBR knowledge soundness of \cref{c:core_iop} }
\end{theorem}
\begin{proof}
\end{proof}

\section{Case-study: SHA-256 hashing}
\albert{Maybe Keccak better?}

\subsection{Arithmetization over $\ZZ[X]$}

\subsubsection{SHA-256 compression overview}
We begin with a brief overview of SHA-256 to establish notation, then formalize the relation we seek to arithmetize.



For the purposes of SHA-256, a \emph{word} is a $\leq 32$ string of bits, which we represent as a binary integer polynomial, i.e.\ as an element of $\{0,1\}^{<32}[X]$.
The SHA-256 compression function receives as input a $512$-bit string, represented as sixteen \emph{message words} $(M_t)_{t\in[16]}\in\{0,1\}^{<32}[X]^{16}$. 
The algorithm also uses $64$ fixed \emph{round constants} $(K_t)_{t\in[64]}\in\{0,1\}^{<32}[X]^{64}$. 

We let 
%
\begin{equation}
\begin{aligned}
\psi_2: \ZZ[X]& \to \ZZ\\
a(X) &\mapsto a(2)
\end{aligned}
\end{equation}
be the ring homomorphism sending integer polynomials to their evaluation at $X=2$. Notice that, given $a\in \{0,1\}^{<32}[X]$, the integer $\psi_2(x) \in \ZZ$ is the integer whose bit-representation coincides with the coefficients of $a$.


SHA-256 is built from the following six auxiliary functions.
For $x,y,z\in\{0,1\}^{<32}[X]$, define:
\begin{equation}\label{eq:sha256_auxiliary_functions}
\begin{aligned}
\Sigma_0(x) &\coloneqq \ROTR^2(x)\xor\ROTR^{13}(x)\xor\ROTR^{22}(x) &&\in\{0,1\}^{<32}[X]\\
\Sigma_1(x)&\coloneqq \ROTR^6(x)\xor\ROTR^{11}(x)\xor\ROTR^{25}(x)&&\in\{0,1\}^{<32}[X]\\
\Ch(x,y,z) &\coloneqq (x\band y)\xor(\Not x\band z)&&\in\{0,1\}^{<32}[X]\\
\Maj(x,y,z)&\coloneqq (x\band y)\xor(x\band z)\xor(y\band z)&&\in\{0,1\}^{<32}[X]\\
\sigma_0(x) &\coloneqq \ROTR^7(x)\xor\ROTR^{18}(x)\xor\SHR^3(x)&&\in\{0,1\}^{<32}[X]\\
\sigma_1(x)&\coloneqq \ROTR^{17}(x)\xor\ROTR^{19}(x)\xor\SHR^{10}(x)&&\in\{0,1\}^{<32}[X]
\end{aligned}
\end{equation}
\albert{todo: define the bit operations}
%The functions $\Ch$ (``choose'') provide $\Maj$ (``majority'') are bitwise selection and voting operations, respectively.

\paragraph{Message schedule.}
At the beginning of the SHA-256 compression function, the sixteen input message words $(M_t)_{t\in[16]}$ are expanded into a so-called \emph{message schedule}, which consists of $64$ integers of $32$ bits $(W_t)_{t\in[64]}\in\ZZ^{64}$ defined as follows:
\begin{equation}\label{eq:sha256_message_schedule_def}
W_t \coloneqq 
\begin{cases}
\psi_2(M_t) & t\in[16],\\
\psi_2(W_{t-16})+\psi_2(\sigma_0(W_{t-15}))+\psi_2(W_{t-7})+\psi_2(\sigma_1(W_{t-2})) \pmod{2^\wordsize} & t\in\{17,\dots,64\}.
\end{cases}
\end{equation}
%The first $16$ schedule words coincide with the input message words; the remaining $48$ are derived by mixing earlier schedule words via the $\sigma$-functions.

\paragraph{Compression loop.}
After the message schedule is computed, SHA-256 performs the so-called \emph{compression loop}. Throughout it, the function maintains an $8$-word internal state, conventionally labeled as $$(a,b,c,d,e,f,g,h) \in \{0,1\}^{<32}[X]^8.$$
The compression loop iterates $64$ rounds.
At each round $t\in[64]$, the state entering that round is denoted
$(a_t,b_t,c_t,d_t,e_t,f_t,g_t,h_t)\in\{0,1\}^{<32}[X]^8$. The per-round update is given by, for $t\in [64]$,
\begin{equation*}
\begin{aligned}
T_{1,t} &\coloneqq \psi_2(h_t)+\psi_2(\Sigma_1(e_t))+\psi_2(\Ch(e_t,f_t,g_t))+\psi_2(K_t)+W_t  \pmod{2^\wordsize} &&\in \ZZ,\\
T_{2,t} &\coloneqq \psi_2(\Sigma_0(a_t))+\psi_2(\Maj(a_t,b_t,c_t))  \pmod{2^\wordsize} &&\in \ZZ,\\
a_{t+1} &\coloneqq \psi_2^{-1}(T_{1,t}+T_{2,t} \pmod{2^\wordsize}) &&\in\{0,1\}^{<32}[X],\\
e_{t+1} &\coloneqq \psi_2^{-1}(\psi_2(d_t)+\psi_2(T_{1,t}) \pmod{2^\wordsize}) &&\in \{0,1\}^{<32}[X],\\
(b_{t+1},&c_{t+1},d_{t+1},f_{t+1},g_{t+1},h_{t+1}) 
\coloneqq (a_t,b_t,c_t,e_t,f_t,g_t)&&\in\{0,1\}^{<32}[X]^6.
\end{aligned}
\end{equation*}
The full SHA-256 compression function adds the initial state to the terminal state word-wise modulo $2^\wordsize$ (the so-called \emph{feed-forward step}); we omit this step for brevity \albert{Let's not do that}. For our purposes, the output of the compression function is the final state $(a_{65},\ldots,h_{65})$.  

\subsubsection{Our arithmetization}

\subsection{End-to-end IOP}

\subsection{Experiments}

\albert{GHASH or b127?}

\albert{Comparison with Flock on SHA256: we commit to 8x less bits. This leads to more than 2x memory reduction. 80\% of prover costs are in the grand product GKR. Hence setting rate=1/4 or smaller has little impact on performance. This leads to smaller proofs. Prover time overall seems to be about 1.5-2.5x worse for us. I also switched to a better config of Ligerito (called k=4 in impl), which reduces proof size by almost 2x, but increases open time, which for us adds a negligible overhead, while in Flock it adds more overhead}

Optimization ideas
\begin{itemize}
    \item b127 instead of GHASH
    \item Lookup table (algorithm-wise, not SNARK-wise) for first levels of GKR
    \item If we have an MLE claim of a vector and its shift by one position, can we reuse something for the latter claim from the computations made when proving the former?
\end{itemize}

\appendix

\section{Ring switching via Galois orbits}\label{s:ring_switch}

\albert{Rewriting the previous appendix for  arbitrary inner products}


%\newcommand{\XX}{\mathbf{X}}

Let $b_1,\ldots, b_m \in \FF_2^{2^n}$.
Using an $\mathbb F_2$-basis $\beta_1,\ldots, \beta_{m}$ of $\mathbb F_{2^m}$ we
pack $m$ inner products %of $2^{m}$ coeffs in $\mathbb F_2$,
\[
\langle b_i, \vv\rangle = \sum_{j\in \bits^{n}} b_{i,j} \cdot \vv_j, 
\quad i = 1, \ldots, m,
\]
coefficient-wise into a single inner product 
\begin{align*}
b&\defeq b_1\cdot \beta_1 + \ldots b_m \cdot \beta_m \in \FF_{2^m}^{2^n}\\
\langle b, \vv\rangle &= \langle b_1, \vv\rangle \cdot \beta_1 + \ldots + \langle b_m, \vv\rangle \cdot\beta_m
\\
&= \sum_{j\in \bits^n} (b_{1,j}\cdot \beta_1 + \ldots + b_{m,j}\cdot \beta_{m}) \cdot \vv_j.
\end{align*}
%
The Frobenius automorphism $\phi(z) = z^2$ keeps $\mathbb F_2\subset \FF_{2^m}$ fixed, and has no other fixed points. 
% (Recall that it is also additive.)
We extract the value of $\langle b_i, \vv\rangle$ from the values of $\langle b, \vv\rangle$ along the Galois orbit 
\[
\vv, \vv^2, \vv^{2^2}, \ldots, \vv^{2^{m-1}},
\]
as follows, where by $\vv^{t}$ we mean that each component of $\vv$ is raised to the power $t$.
%
The Moore matrix
\begin{align*}
M = \begin{pmatrix}
\beta_1 & \beta_2 & \ldots & \beta_m
\\
\beta_1^2 & \beta_2^2 & \ldots & \beta_m^2
\\
\beta_1^{2^2} & \beta_2^{2^2} & \ldots & \beta_m^{2^2}
\\
\vdots
\\
\beta_1^{2^{m-1}} & \beta_2^{2^{m-1}} & \ldots & \beta_m^{2^{m-1}}
\end{pmatrix}
\end{align*}
consists of the Galois orbits of the basis vectors, and is invertible. (See below.)
Using that the component polynomials are over $\mathbb F_2$, we see that
%We write $v_{k} = b(\phi^{-k} z)$, $k=0, \ldots, m-1$, and 
% \[
% v_k = \phi^{-k}\left(b_1(z) \cdot \phi^{k}\beta_1 + \ldots + b_m(z) \cdot \phi^{k}\beta_m \right),
% \]
% since $b_i(\phi^{-k} z) = \phi^{-k} b_i(z)$, in other words
\[
\phi^{k} \left(\langle b, \phi^{-k}(\vv)\rangle\right) =  \langle b_1, \vv\rangle  \cdot \phi^{k}\beta_1 + \ldots + \langle b_m, \vv\rangle   \cdot \phi^{k}\beta_m,
\]
for any $k$. 
%
To this end, observe that
\[
M \cdot\begin{pmatrix}
\langle b_1, \vv\rangle  
\\
\langle b_2, \vv\rangle  
\\
\langle b_3, \vv\rangle  
\\
\vdots
\\
\langle b_m, \vv\rangle  
\end{pmatrix}
=
\begin{pmatrix}
\langle b, \vv\rangle
\\
\phi \left(\langle b, \phi^{-1}(\vv)\rangle\right)
\\
\phi^{2} \left(\langle b, \phi^{-2}(\vv)\rangle\right)
\\
\vdots
\\
\phi^{m-1} \left(\langle b, \phi^{-(m-1)}(\vv)\rangle\right)
\end{pmatrix}.
\]
Multiplying both sides by $M^{-1}$ we obtain an expression of the values $\langle b_i, \vv\rangle$ as a linear expression of the values $\phi^{k} \left(\langle b, \phi^{-k}(\vv)\rangle\right)$.
% and $M\cdot c$ is the Galois orbit of $c = (c_1,\ldots, c_m)$ in coordinates of the basis.
% The rows of $M^{-1}$ give us the linear functionals that recover the $i$-th coordinate of $z$ from its Galois-orbit.
\ulrich{%
\small
Invertibility of $M$ follows from that of $M^t$: 
any linear functional $(c_1,\ldots, c_m)$ in the kernel of $M^t$ gives rise to the linearizable polynomial
\[
p(X) = \sum_{i=0}^{m-1} c_i\cdot X^{2^i} 
\]
which has $\beta_1,\ldots, \beta_m$, and hence the whole space $F$ as roots. 
By dimension, $p(X)$ must be trivial, meaning that all its coefficents are zero.
}






\section{[Old appendix by Ulrich] Galois FRI-Binius}

\ulrich{Just notes, not sure if it worth to keep it in final doc} \albert{Nice, thanks!}

Packed commitment scheme for $\mathbb F_{2^m}$ . % in the values-to-uni-coeffs regime.

For simplicity we assume packing across polynomials. 
(The single poly case goes likewise.)
Using an $\mathbb F_2$-basis $\beta_1,\ldots, \beta_{m}$ of $\mathbb F_{2^m}$ we
pack $m$ polys of $2^n$ coefficients %of $2^{m}$ coeffs in $\mathbb F_2$,
\[
b_i(X) = \sum_{j=0}^{2^n - 1} b_{i,j} \cdot X^i \quad \in \mathbb F_2[X]^{< 2^n}, 
\quad i = 1, \ldots, m,
\]
coefficient-wise into a single poly 
\begin{align*}
b(X) &= b_1(X)\cdot \beta_1 + \ldots + b_{m}(X) \cdot\beta_m
\\
&= \sum_{j=0}^{2^n} (b_{1,j}\cdot \beta_1 + \ldots + b_{m,j}\cdot \beta_{m}) \cdot X^j 
\quad \in \mathbb F_{2^m}[X]^{<2^n}.
\end{align*}

The Frobenius automorphism $\phi(z) = z^2$ keeps $\mathbb F_2\subset F$ fixed, and has no other fixed points. 
% (Recall that it is also additive.)
For any point $z\in \mathbb F$, we extract the value of $b_i(z)$ from the values of the packed poly $b(X)$ along the Galois orbit 
\[
z, z^2, z^{2^2}, \ldots, z^{2^{m-1}}
\]
as follows.
%
The Moore matrix
\begin{align*}
M = \begin{pmatrix}
\beta_1 & \beta_2 & \ldots & \beta_m
\\
\beta_1^2 & \beta_2^2 & \ldots & \beta_m^2
\\
\beta_1^{2^2} & \beta_2^{2^2} & \ldots & \beta_m^{2^2}
\\
\vdots
\\
\beta_1^{2^{m-1}} & \beta_2^{2^{m-1}} & \ldots & \beta_m^{2^{m-1}}
\end{pmatrix}
\end{align*}
consists of the Galois orbits of the basis vectors, and is invertible. (See below.)
Using that the component polynomials are over $\mathbb F_2$, we see that
%We write $v_{k} = b(\phi^{-k} z)$, $k=0, \ldots, m-1$, and 
% \[
% v_k = \phi^{-k}\left(b_1(z) \cdot \phi^{k}\beta_1 + \ldots + b_m(z) \cdot \phi^{k}\beta_m \right),
% \]
% since $b_i(\phi^{-k} z) = \phi^{-k} b_i(z)$, in other words
\[
\phi^{k} \left(b(\phi^{-k} z)\right) = b_1(z) \cdot \phi^{k}\beta_1 + \ldots + b_m(z) \cdot \phi^{k}\beta_m,
\]
for any $k$. 
To extract the $b_1(z)$ term, we look for weights $c_1,\ldots, c_m\in F$ which annihilate the orbits of the other $\beta_2,\ldots, \beta_m$, concretely
\begin{align*}
\sum_{k=0}^{m-1} c_k \cdot \phi^k \beta_1 = 1,
\quad 
\sum_{k=0}^{m-1} c_k \cdot \phi^k \beta_j = 0 \quad j=2,\ldots,n.
\end{align*}
This is exactly the first row vector of $M^{-1}$. 
%
In general, 
\[
\begin{pmatrix}
b_1(z)
\\
b_2(z)
\\
b_3(z)
\\
\vdots
\\
b_m(z)
\end{pmatrix}
= M^{-1}\cdot 
\begin{pmatrix}
v_0
\\
\phi v_1
\\
\phi^2 v_2
\\
\vdots
\\
\phi^{m-1} v_{m-1}
\end{pmatrix},
\]
where $v_k = b(\phi^{-k} z)$.
% and $M\cdot c$ is the Galois orbit of $c = (c_1,\ldots, c_m)$ in coordinates of the basis.
% The rows of $M^{-1}$ give us the linear functionals that recover the $i$-th coordinate of $z$ from its Galois-orbit.
\ulrich{%
\small
Invertibility of $M$ follows from that of $M^t$: 
any linear functional $(c_1,\ldots, c_m)$ in the kernel of $M^t$ gives rise to the linearizable polynomial
\[
p(X) = \sum_{i=0}^{m-1} c_i\cdot X^{2^i} 
\]
which has $\beta_1,\ldots, \beta_m$, and hence the whole space $F$ as roots. 
By dimension, $p(X)$ must be trivial, meaning that all its coefficents are zero.
}




\section{Bilinear Embeddings}

Given a field $\mathbb F$ and a degree-$d$ extension $\mathbb E$. Our goal is to commit to a vector $\mathbf{w} \in \mathbb{F}^{n d}$ so that we can later use it to prove an inner product claim of the form $\langle \mathbf{w}, \mathbf{a}\rangle_\mathbb{F} = h$ for some $\mathbf{a} \in \mathbb{F}^{n d}$ and $h \in \mathbb{F}$. However, we only have an efficient inner product commitment scheme to commit to vectors $\widetilde{\mathbf{w}} \in \mathbb{E}^{n}$ and open as $\langle \widetilde{\mathbf{w}}, \widetilde{\mathbf{a}}\rangle_\mathbb{E} = \widetilde{h}$ for some $\widetilde{\mathbf{a}} \in \mathbb{E}^{n}$ and $\widetilde{h} \in \mathbb{E}$.

To do this, we seek a triple of maps $(W, A, H)$ that satisfy:
\[
\langle \mathbf{w},\mathbf{a}\rangle_{\mathbb{F}}
=
H\left(\langle W(\mathbf{w}), A(\mathbf{a}) \rangle_{\mathbb{E}} \right).
\]

\subsection{Background}
In the following let $\mathbb{E}/\mathbb{F}$ be a finite field extension of degree $d$. If we disregard multiplication in $\mathbb{E}$, then $\mathbb{E}$ is a $d$-dimensional vector space over $\mathbb{F}$, hence isomorphic to $\mathbb{F}^d$. Let $\mathbf{e}_1,\ldots,\mathbf{e}_d$ denote the standard basis of $\mathbb{F}^d$.
Recall that a function $\mu:\mathbb{E}\times \mathbb{E} \rightarrow \mathbb{F}$ is a \emph{bilinear form}  if it is linear in each argument separately. It is \emph{symmetric} if $\mu(x, y) = \mu(y, x)$ for all $x,y \in \mathbb{E}$.
For each $x \in \mathbb{E}$, the partial application in the first argument defines a linear form
\[
\mu_{x} : \mathbb{E} \rightarrow \mathbb{F} \quad y \mapsto \mu(x, y).
\]
Thus $\mu_{x} \in \mathbb{E}^*$ where $\mathbb{E}^* = \operatorname{Hom}_\mathbb{F}(\mathbb{E}, \mathbb{F})$ denotes the dual vector space. The bilinear form therefore induces an $\mathbb{F}$-linear map
\[
\Phi_\mu:\mathbb{E}\to\mathbb{E}^* \qquad x\mapsto\mu_x.
\]
We call $\mu$ \emph{nondegenerate}, or equivalently a \emph{perfect pairing}, if $\Phi_\mu$ is an isomorphism. Or equivalently, \albert{add a comment that injectivity implies surjectivity when the vector spaces have finite dimension?}
\[
\mu(x,y)=0\text{ for all }y\in\mathbb E \quad\Longrightarrow\quad x=0.
\]
Since $\mu$ is a perfect pairing, every basis $(\beta_1,\ldots,\beta_d)$ of $\mathbb{E}$ has a unique dual basis $(\beta_1^\vee,\ldots,\beta_d^\vee)$ satisfying $\mu(\beta_i, \beta_j^\vee) = \delta_{ij}$.

\subsection{Structure}
We are now equipped to prove a theorem on the structure of embeddings. Under the assumption that the output map is $\mathbb{F}$-linear, we show that for $n=1$ the space of possible embeddings is very structured. This then extends naturally to $n>1$ by doing blockwise packing.

\begin{theorem}[Structure of bilinear embeddings]
Let $W,A:\mathbb{F}^d\rightarrow\mathbb{E}, H:\mathbb{E}\rightarrow\mathbb{F}$ be maps such that for all
$\mathbf{w},\mathbf{a}\in\mathbb{F}^d$,
\[
\langle \mathbf{w},\mathbf{a}\rangle_{\mathbb{F}}
=
H\bigl(W(\mathbf{w})A(\mathbf{a})\bigr).
\]
Suppose that $H$ is $\mathbb{F}$-linear. Then: \albert{It would be nice to also add that $A$ and $W$ are bijective, in particular they are the namesake ``embedings''}
\begin{enumerate}
    \item $W$ and $A$ are $\mathbb{F}$-linear,
    \item $\mu_H(x,y) = H(xy)$ is a perfect pairing,
    \item $\left(W(\mathbf{e}_1),\ldots,W(\mathbf{e}_d)\right)$ and $\left(A(\mathbf{e}_1),\ldots,A(\mathbf{e}_d)\right)$ are dual bases with respect to $\mu_H$.
\end{enumerate}

Conversely, if $W$ and $A$, and $H$ are $\mathbb{F}$-linear, $\mu_H$ is a perfect pairing and the images of the standard basis under $W$ and $A$ are dual bases with respect to $\mu_H$, then the identity
\[
\langle \mathbf{w},\mathbf{a}\rangle_{\mathbb{F}}
=
H\bigl(W(\mathbf{w})A(\mathbf{a})\bigr).
\]
holds for all $\mathbf{w},\mathbf{a}\in\mathbb{F}^d$.
\end{theorem}

\begin{proof}[Proof of the forward direction]
First, observe that $
\mu_H(W(\mathbf e_i),A(\mathbf e_j))
= \langle \mathbf e_i,\mathbf e_j\rangle_{\mathbb F}
= \delta_{ij}
$. Taking $i=j$ shows $H$ is nonzero. Since $H$ is $\mathbb{F}$-linear, $\mu_H$ is a symmetric bilinear form, we show that it is nondegenerate. Let $x\in\mathbb E$ be nonzero. Since multiplication by $x$ is a bijection of $\mathbb E$ and $H$ is nonzero, there exists $y\in\mathbb E$ such that $\mu_H(x,y)=H(xy)\neq 0$. Thus $\mu_H$ is nondegenerate, and hence a perfect pairing.

Next, we show that $W$ and $A$ are bijective. Suppose
$W(\mathbf w_1)=W(\mathbf w_2)$. Then for every
$\mathbf a\in\mathbb F^d$,
\[
\langle \mathbf w_1-\mathbf w_2,\mathbf a\rangle_{\mathbb F}
= \mu_H(W(\mathbf w_1) - W(\mathbf w_2), A(\mathbf a))
=0.
\]
By nondegeneracy of the standard inner product on $\mathbb F^d$, this
implies $\mathbf w_1=\mathbf w_2$. Hence $W$ is injective. Since
$\mathbb F^d$ and $\mathbb E$ have the same finite cardinality, $W$ is
bijective. The same argument shows that $A$ is bijective.

We now prove that $W$ is $\mathbb F$-linear. For
$\mathbf w_1,\mathbf w_2,\mathbf a\in\mathbb F^d$, linearity of $H$
gives
\[
\begin{aligned}
&\mu_H\bigl(
W(\mathbf w_1+\mathbf w_2)-W(\mathbf w_1)-W(\mathbf w_2),
A(\mathbf a)
\bigr)\\
&=
\langle \mathbf w_1+\mathbf w_2,\mathbf a\rangle_{\mathbb F}
-
\langle \mathbf w_1,\mathbf a\rangle_{\mathbb F}
-
\langle \mathbf w_2,\mathbf a\rangle_{\mathbb F}\\
&=0.
\end{aligned}
\]
Since $A$ is surjective, the second argument ranges over all of
$\mathbb E$. Nondegeneracy of $\mu_H$ therefore implies
$W(\mathbf w_1+\mathbf w_2) = W(\mathbf w_1)+W(\mathbf w_2)$. Similarly, for $\lambda\in\mathbb F$,
\[
\begin{aligned}
&\mu_H\bigl(
W(\lambda\mathbf w)-\lambda W(\mathbf w),
A(\mathbf a)
\bigr)\\
&=
\langle \lambda\mathbf w,\mathbf a\rangle_{\mathbb F}
-
\lambda\langle \mathbf w,\mathbf a\rangle_{\mathbb F}\\
&=0.
\end{aligned}
\]
Again using surjectivity of $A$ and nondegeneracy of $\mu_H$, we obtain $W(\lambda\mathbf w)=\lambda W(\mathbf w)$. Hence $W$ is $\mathbb F$-linear. By symmetry, $A$ is also
$\mathbb F$-linear. Since they are also bijective, $W$ and $A$ are linear isomorphisms. From this it follows that $\left(W(\mathbf e_1),\ldots,W(\mathbf e_d)\right)$ and $
\left(A(\mathbf e_1),\ldots,A(\mathbf e_d)\right)$ are bases of $\mathbb E$. Finally, from $\mu_H(W(\mathbf e_i),A(\mathbf e_j)) = \delta_{ij}$ it follows they are dual with respect to $\mu_H$.
\end{proof}

\begin{proof}[Proof of the converse direction.] Can be directly shown using a basis expansion. 
Write $\mathbf w=\sum_{i=1}^d w_i\mathbf e_i$ and $\mathbf a=\sum_{j=1}^d a_j\mathbf e_j$
Since $W$ and $A$ are $\mathbb F$-linear, $W(\mathbf w) = \sum_{i=1}^d w_iW(\mathbf e_i)$ and $A(\mathbf a) = \sum_{j=1}^d a_jA(\mathbf e_j)$. Therefore,
\[
\begin{aligned}
H\bigl(W(\mathbf w)A(\mathbf a)\bigr)
&=
\mu_H\bigl(W(\mathbf w),A(\mathbf a)\bigr)\\
&=
\sum_{i,j}
w_i a_j
\mu_H\bigl(W(\mathbf e_i),A(\mathbf e_j)\bigr)\\
&=
\sum_{i,j}
w_i a_j\delta_{ij}\\
&=
\sum_iw_ia_i\\
&=
\langle\mathbf w,\mathbf a\rangle_{\mathbb F}.
\end{aligned}
\]
\end{proof}

We have now shown that every embedding with a linear output map is determined by a choice of a nonzero linear map $H$ and linear isomorphism $W$. The map $W$ then determines the first basis $\beta_i=W(\mathbf e_i)$ and the map $A$ is then forced to map to the $\mu_H$-dual basis.

The linearity assumption on the output map is necessary. Indeed, composing a linear embedding with a multiplicative bijection $g:\mathbb E\to\mathbb E$ generally produces an embedding that is nonlinear. Whether every nonlinear embedding arises from such a multiplicative reparameterization is a separate question, which we do not need here.

The set of nonzero linear output maps $H$ is $\mathbb{E}^* \setminus \{0\}$ which is isomorphic to $\mathbb{E}^\times$. The set of linear isomorphism $W$ is isomorphic to $\mathrm{GL}_d(\mathbb{F})$, hence the total number of bilinear embeddings is $|\mathbb{E}^\times| \cdot |\mathrm{GL}_d(\mathbb{F})|$. Let's now look at specific choices.

\subsection{Implementation}

\ulrich{Let us copy paste this into the new f2z doc. You are working on the legacy overleaf doc.}
\begin{theorem}[Preservation of MLE openings] Given $d = 2^k$, $n = 2^l$ and a $\mathbb{F}$-linear map $A : \mathbb{F}^d \to \mathbb{E}$ applied blockwise, then if $\mathbf{a}\in \mathbb{F}^{2^{k + l}}$ given by $a_i = \mathrm{eq}_\mathbb{F}(\mathbf{r}, \mathbf{b}(i))$ where $\mathbf{b}$ is the binary encoding, the embedding $A(\mathbf{a}) \in \mathbb{E}^{2^l}$ is given by
\[
(A(\mathbf{a}))_i = \alpha \cdot \mathrm{eq}_\mathbb{E}(\mathbf{r}', \mathbf{b}(i))
\]
for some $\alpha \in \mathbb{E}$ and $\mathbf{r}' \in \mathbb{E}^l$.
\end{theorem}

In particular this means that any bilinear map can also be used to lift an multilinear commitment scheme.

\begin{theorem}[Preservation of succinct verification] If $\mathbf{a}$ has a Matrix-branching program for succinct evaluation of its multilinear extension, then $A(\mathbf{a})$ also has one, at most $d$ times wider. 
\end{theorem}


\subsection{Trace embedding}

Let $q = |\mathbb F|$, the $q$-power Frobenius map is $\varphi:\mathbb E\to\mathbb E$ given by $x\mapsto x^q$. It is an $\mathbb F$-linear field automorphism of order $d$, and its
fixed field is precisely $\mathbb F$:
\[
\mathbb E^\varphi
= \{x\in\mathbb E:\varphi(x)=x\}
= \mathbb F.
\]
The field trace is the $\mathbb F$-linear map $\operatorname{Tr}_{\mathbb E/\mathbb F}:\mathbb E\to\mathbb F$ defined by
\[
\operatorname{Tr}_{\mathbb E/\mathbb F}(x)
= \sum_{i=0}^{d-1}\varphi^i(x)
= \sum_{i=0}^{d-1}x^{q^i}.
\]
To see that this sum is indeed in $\mathbb{F} \subseteq \mathbb{E}$ observe the sum is fixed by $\varphi$, since
\[
\varphi\left( \sum_{i=0}^{d-1}\varphi^i(x) \right)
= \sum_{i=1}^{d}\varphi^i(x)
= \sum_{i=0}^{d-1}\varphi^i(x),
\]
where $\varphi^d(x)=\varphi^0(x)$. Hence the trace lies in the fixed field $\mathbb F$.

This makes $\operatorname{Tr}_{\mathbb E/\mathbb F}$ a natural choice of output map $H$. More generally, every $\mathbb F$-linear map $H:\mathbb E\to\mathbb F$ is uniquely of the form $H(x)=\operatorname{Tr}_{\mathbb E/\mathbb F}(\alpha x)$ for some $\alpha\in\mathbb E$. Consequently, the nonzero output maps are parameterized by $\mathbb E^\times$.

The Galois orbit embedding chooses the basis $\beta_i=\varphi^i(\gamma)$ which has dual $\beta_i^\vee=\varphi^i(\gamma^\vee)$.

\ulrich{%
This sounds to me very much related to Galois Fribinius on an inner product with a general covector. 
You would end up with the multi-claim of the inner products of the packed poly with the Galois orbit of the covector.
This is how I view it:
Given $\vec p_1, \ldots, \vec p_d \in \mathbb F^n$, and $\vec q_1, \ldots, \vec q_d \in \mathbb F^n$, we are interested in 
\[
\sum_{k=1}^d \langle \vec p_k, \vec q_k\rangle,
\]
right?

The packing map sends the coeff vectors $\vec p_1,\ldots, \vec p_d$ to $\vec p = \beta_1\cdot \vec p_1 + \ldots + \beta_d\cdot \vec p_d$, where $\beta_1,\ldots,\beta_d$ is an $\mathbb F$-basis of $\mathbb E$. 
Choose the element $a_k$ in $\mathbb E$ which extracts the $k$-th coordinate
\[
\vec p_k = \Phi_k(\vec p) = \operatorname{Tr}_{\mathbb E/\mathbb F}(a_k\cdot \vec p).
\]
Then 
\begin{align*}
\langle \vec p_k, \vec q_k\rangle 
&= \langle \sum_{i=0}^d \varphi^i(a_k\cdot \vec p), \vec q_k\rangle 
\\
&= \sum_{i=0}^d  \varphi^{i}(a_k) \cdot \varphi^i(\langle \vec p,  \vec q_k\rangle),
\end{align*}
since $\vec q_k$ is over $\mathbb F$.
This reduces our inner product claim to that of the packed poly $\vec p$ with the unpacked polys $\vec q_1,\ldots, \vec q_k$.

But does that match our use case?
We are rather interested in $\langle \vec p_k, \vec v\rangle$ for some $\vec v\in \mathbb E^n$.
In this case, 
\begin{align*}
\langle \vec p_k, \vec v \rangle 
&= \langle \sum_{i=0}^d \varphi^i(a_k\cdot \vec p), \vec q_k\rangle 
\\
&= \sum_{i=0}^d  \varphi^{i}(a_k) \cdot \varphi^i(\langle \vec p,  \varphi^{-i}(\vec v)\rangle),
\end{align*}
and we end up with inner products of the packed poly with the Galois orbit of $\vec v$. 
}


\subsection{Coefficient projection}

In the coefficient projection we assume the standard polynomial representation for $\mathbb{E}$ (as is common in practical implementations) and take $H: \mathbb{E} \to \mathbb{F}$ to select a particular coordinate, e.g. the constant coordinate. We will show that this lead to a particularly efficient implementation where $W$ is the identity map and $A$ is generally sparse and for many fields monomial.


\end{document}
